\subsection{谱与预解式}

引理 4. --- 设 E 为复巴拿赫空间，u 为 E 上的闭单射部分算子，满足 \(u^{-1} \in \mathcal{L}(\mathrm{E})\)。设 \(v \in \mathcal{L}(\mathrm{E})\)

满足 \(\|v\| < \|u^{-1}\|^{-1}\)。则部分算子 \(u + v\) 是单射的，我们有 \((u + v)^{-1} \in \mathcal{L}(\mathrm{E})\)，并且

\[
(u + v) ^ {- 1} = u ^ {- 1} \circ \sum_ {k = 0} ^ {+ \infty} (- v u ^ {- 1}) ^ {k},\tag{6}
\]

其中该级数在 \(\mathcal{L}(\mathrm{E})\) 中绝对收敛。此外，我们有

\[
\| (u + v) ^ {- 1} \| \leqslant \frac {\| u ^ {- 1} \|}{1 - \| v \| \| u ^ {- 1} \|}.
\]

由于 \(\|v\|\|u^{-1}\|<1\)，通项为 \((-vu^{-1})^{k}\) 的级数在 \(\mathcal{L}(\mathrm{E})\) 中绝对收敛，且其和是自同态 \(1_{E}+vu^{-1}\) 的逆（I，第 22 页，命题 2）。于是，部分算子 \((1_{\mathrm{E}}+vu^{-1})\circ u=u+v\)（IV，第 226 页，注 vi）是单射的，且逆部分算子 \((u+v)^{-1}=u^{-1}\circ(1_{\mathrm{E}}+vu^{-1})^{-1}\)（IV，第 226 页，注 iii）属于 \(\mathcal{L}(\mathrm{E})\)。由于

\[
\left\| \left(1 _ {\mathrm{E}} + v u ^ {- 1}\right) ^ {- 1} \right\| \leqslant \sum_ {k = 0} ^ {+ \infty} (\left\| v \right\| \left\| u ^ {- 1} \right\|) ^ {k} = \frac {1}{1 - \left\| v \right\| \left\| u ^ {- 1} \right\|},
\]

引理得证。

定义 8. --- 设 u 是复巴拿赫空间 E 上的稠定闭算子。u 的预解集是由这样的复数 λ 组成的集合：部分算子 \(\lambda 1_{\mathrm{E}} - u\) 是单射的，且其逆 \((\lambda 1_{\mathrm{E}} - u)^{-1}\) 属于 \(\mathcal{L}(\mathrm{E})\)。

u 的谱，记作 \(\mathrm{Sp}(u)\)，是预解集在 C 中的余集。

若 \(\lambda \in \mathbf{C} - \mathrm{Sp}(u)\)，我们用 \(\mathrm{R}(u,\lambda) \in \mathcal{L}(\mathrm{E})\) 表示 \(\lambda 1_{\mathrm{E}} - u\) 的逆。从 \(\mathbf{C} - \mathrm{Sp}(u)\) 到 \(\mathcal{L}(\mathrm{E})\) 中把 \(\lambda\) 映为 \(\mathrm{R}(u,\lambda)\) 的映射称为 \(u\) 的预解式。

注记。--- 设 E 为复巴拿赫空间，u 为 E 上稠定的闭算子。

\begin{enumerate}
\def\labelenumi{\arabic{enumi})}
\item
  若 \(u \in \mathcal{L}(\mathrm{E})\)，则它的谱与元素 \(u\) 作为代数 \(\mathcal{L}(\mathrm{E})\) 中元素的谱一致（I，第 2 页，定义 1）。
\item
  设 D 为 \(u\) 的定义域。对每个 \(\lambda \in \mathbf{C} - \mathrm{Sp}(u)\)，我们有
\end{enumerate}

\[
1 _ {\mathrm{E}} = (\lambda 1 _ {\mathrm{E}} - u) \circ \mathrm{R} (u, \lambda), \quad 1 _ {\mathrm{D}} = \mathrm{R} (u, \lambda) \circ (\lambda 1 _ {\mathrm{E}} - u).\tag{7}
\]

此外，对于 u 的预解集中的 \(\lambda_{1}\) 与 \(\lambda_{2}\)，我们有

\[
\mathrm{R} (u, \lambda_ {1}) - \mathrm{R} (u, \lambda_ {2}) = (\lambda_ {2} - \lambda_ {1}) \mathrm{R} (u, \lambda_ {2}) \circ \mathrm{R} (u, \lambda_ {1}),\tag{8}
\]

\[
\mathrm{R} (u, \lambda_ {1}) \circ \mathrm{R} (u, \lambda_ {2}) = \mathrm{R} (u, \lambda_ {2}) \circ \mathrm{R} (u, \lambda_ {1}),\tag{9}
\]

事实上，有

\[
\mathrm{R} (u, \lambda_ {1}) - \mathrm{R} (u, \lambda_ {2}) = \mathrm{R} (u, \lambda_ {1}) \circ 1 _ {\mathrm{E}} - 1 _ {\mathrm{D}} \circ \mathrm{R} (u, \lambda_ {2}).
\]

由于 \(1_{\mathrm{E}} = (\lambda_21_{\mathrm{E}} - u)\circ \mathrm{R}(u,\lambda_2)\) 与 \(1_{\mathrm{D}} = \mathrm{R}(u,\lambda_1)\circ (\lambda_11_{\mathrm{E}} - u)\)，我们得到

\[
\begin{array}{c} \mathrm{R} (u, \lambda_ {1}) - \mathrm{R} (u, \lambda_ {2}) = \lambda_ {2} \mathrm{R} (u, \lambda_ {1}) \circ \mathrm{R} (u, \lambda_ {2}) \\ - \mathrm{R} (u, \lambda_ {1}) \circ u \circ \mathrm{R} (u, \lambda_ {2}) - \lambda_ {1} \mathrm{R} (u, \lambda_ {1}) \circ \mathrm{R} (u, \lambda_ {2}) \\ + \mathrm{R} (u, \lambda_ {1}) \circ u \circ \mathrm{R} (u, \lambda_ {2}), \end{array}
\]

由此得到第一个公式。交换 \(\lambda_{1}\) 与 \(\lambda_{2}\) 的角色，便推出第二个公式。

\begin{enumerate}
\def\labelenumi{\arabic{enumi})}
\setcounter{enumi}{2}
\item
  设 \(\lambda \in \mathbf{C}\)。由闭图像定理（EVT，I，第 19 页，推论 5），若从 \(\mathrm{dom}(u)\) 到 E 中的由 \(x \mapsto (\lambda 1_{\mathrm{E}} - u)(x)\) 定义的线性映射是双射，则它的逆（其图像是闭的）从 E 到 E 连续。因此，\(\lambda\) 属于 \(u\) 的预解集，当且仅当由 \(x \mapsto (\lambda 1_{\mathrm{E}} - u)(x)\) 给出的从 \(\mathrm{dom}(u)\) 到 E 中的映射是双射。
\item
  若 \(\lambda\) 属于 u 的谱，则下列性质之一成立：
\end{enumerate}

\begin{enumerate}
\def\labelenumi{(\roman{enumi})}
\item
  \(\lambda 1_{\mathrm{E}} - u\) 的核不化为 0；这时称 \(\lambda\) 是 \(u\) 的特征值，并称 \(\operatorname{Ker}(\lambda 1_{\mathrm{E}} - u)\) 的维数为其重数；
\item
  部分算子 \(\lambda 1_{\mathrm{E}} - u\) 是单射的，而其像在 E 中不稠密；这时称 \(\lambda\) 属于 \(u\) 的剩余谱；
\item
  部分算子 \(\lambda 1_{\mathrm{E}} - u\) 是单射的，其像在 E 中稠密，但 \(\lambda 1_{\mathrm{E}} - u\) 不是满射；这时称 \(\lambda\) 属于 \(u\) 的连续谱。
\end{enumerate}

\begin{enumerate}
\def\labelenumi{\arabic{enumi})}
\setcounter{enumi}{4}
\tightlist
\item
  设 \(\lambda\) 是属于 \(u\) 的预解集的复数。预解式 \(\mathrm{R}(u,\lambda)\) 是从 E 到 E 中的单射线性映射；其像是 \(u\) 的定义域，且 \(u = \lambda 1_{\mathrm{E}} - \mathrm{R}(u,\lambda)^{-1}\)（参见 IV，第 226 页）。反之，这一性质刻画了预解集与预解式。准确地说，设 \(\lambda \in \mathbf{C}\)；若存在从 E 到 E 中的连续单射线性映射 w，使得 \(u = \lambda 1_{E} - w^{-1}\)，则 \(\lambda\) 属于 u 的预解集，且 \(w = \mathrm{R}(u, \lambda)\)。
\end{enumerate}

特别地，若 v 是 E 上的闭稠定算子，且 \(\lambda \in C\) 是不属于 \(\operatorname{Sp}(u) \cup \operatorname{Sp}(v)\) 的复数，则等式 \(\operatorname{R}(u, \lambda) = \operatorname{R}(v, \lambda)\) 蕴含 u = v。

\begin{enumerate}
\def\labelenumi{\arabic{enumi})}
\setcounter{enumi}{5}
\tightlist
\item
  我们把可闭部分算子 \(u\) 的谱定义为其闭包的谱。
\end{enumerate}

存在谱为空集、或谱等于 C 的闭算子（IV，第 347 页，习题 12）。

设 E 为复巴拿赫空间，u 为 E 上的闭稠定算子。若 F 是复巴拿赫空间，且 \(v: E \to F\) 是同构，则我们有 \(\mathrm{Sp}(v \circ u \circ v^{-1}) = \mathrm{Sp}(u)\) 与 \(\mathrm{R}(v \circ u \circ v^{-1}, \lambda) = v \circ \mathrm{R}(u, \lambda) \circ v^{-1}\)，对一切 \(\lambda \notin \mathrm{Sp}(u)\) 成立。

命题 14. --- 设 E 为复巴拿赫空间。设 u 为 E 上稠定的闭算子，U 为其预解集。a) 对每个 \(\lambda\in\mathrm{U}\)，C 中以 \(\lambda\) 为心、以 \(\|\mathrm{R}(u,\lambda)\|^{-1}\) 为半径的开圆盘都包含于 U；

\begin{enumerate}
\def\labelenumi{\alph{enumi})}
\setcounter{enumi}{1}
\item
  集合 U 是 C 中的开集；
\item
  设 \(\operatorname{Sp}(u)\) 非空。设 \(\lambda \in \mathrm{U}\)，\(\delta\) 是 \(\mathbf{C}\) 中从 \(\lambda\) 到 \(u\) 的谱的距离。我们有 \(\delta > 0\) 及 \(\| \mathrm{R}(u, \lambda) \| \geqslant 1 / \delta\)；
\item
  映射 \(\lambda \mapsto \mathrm{R}(u,\lambda)\) 是从 U 到 \(\mathcal{L}(\mathrm{E})\) 中的全纯映射。对每个整数 \(k\in \mathbf{N}\) 与每个 \(\lambda \in \mathrm{U}\)，我们有
\end{enumerate}

\[
\frac {\partial^ {k}}{\partial \lambda^ {k}} \mathrm{R} (u, \lambda) = (- 1) ^ {k} k! \mathrm{R} (u, \lambda) ^ {k + 1}.
\]

设 \(\lambda \in \mathrm{U}\)。对每个 \(\mu \in \mathbf{C}\)，只要 \(\| (\mu -\lambda)1_{\mathrm{E}}\| < \| \mathrm{R}(u,\lambda)\|^{-1}\)，将引理 4 应用于单射部分算子 \(\lambda 1_{\mathrm{E}} - u\) 并取 \(v = (\mu -\lambda)1_{\mathrm{E}}\)，便知部分算子 \(\mu 1_{\mathrm{E}} - u = \lambda 1_{\mathrm{E}} - u + v\) 是单射的且具有连续的逆。这就蕴含 \(a\)）。由同上所引，我们还有公式

\[
\mathrm{R} (u, \mu) = \mathrm{R} (u, \lambda) \circ \sum_ {k \in \mathbf {N}} (\lambda - \mu) ^ {k} \mathrm{R} (u, \lambda) ^ {k},
\]

由此推出 u 的预解式在 U 中全纯。

断言 \(b\)）由 \(a\)）立得。若 \(\mathrm{Sp}(u)\) 非空，则 \(\lambda\) 到 \(\mathrm{Sp}(u)\) 的距离严格为正（TG，IX，第 13 页，命题 2），由此得 \(c\)）。

断言 \(d\)）的最后一部分对 \(k\) 用归纳法证明，其中 \(k = 1\) 的情形是第 245 页公式 (8) 的推论。

命题 15. --- 设 u 是复巴拿赫空间 E 上稠定的闭算子，λ 是属于 u 的预解集的复数。

\begin{enumerate}
\def\labelenumi{\alph{enumi})}
\item
  子集 \(\mathrm{Sp}(\mathrm{R}(u,\lambda)) \setminus \{0\}\)（\(\mathbf{C}\) 中的）是 \(u\) 的谱经映射 \(\mu \mapsto (\lambda -\mu)^{-1}\)（从 \(\mathbf{C} \setminus \{\lambda \}\) 到 \(\mathbf{C}^*\) 中）所成的像；
\item
  对每个 \(\mu \neq \lambda\)（在 \(\mathbf{C}\) 中），有
\end{enumerate}

\[
\operatorname{Ker} \left(\mu 1 _ {\mathrm{E}} - u\right) = \operatorname{Ker} \left(\left(\lambda - \mu\right) ^ {- 1} 1 _ {\mathrm{E}} - \mathrm{R} (u, \lambda)\right).
\]

我们来证明断言 \(a\)）。对每个 \(\mu \neq \lambda\)，计算

\[
\mu 1 _ {\mathrm{E}} - u = (\lambda - \mu) ((\lambda - \mu) ^ {- 1} 1 _ {\mathrm{E}} - \mathrm{R} (u, \lambda)) (\lambda 1 _ {\mathrm{E}} - u).
\]

由于 \(\lambda \notin \mathrm{Sp}(u)\) 且 \(\mu \neq \lambda\)，线性映射 \((\lambda - \mu)(\lambda 1_{\mathrm{E}} - u)\) 是从 \(\mathrm{dom}(u)\) 到 E 上的双射。于是，该公式表明：\(\mu 1_{\mathrm{E}} - u\) 是从 \(\mathrm{dom}(u)\) 到 E 上的双射，当且仅当 \((\lambda - \mu)^{-1} 1_{\mathrm{E}} - \mathrm{R}(u, \lambda)\) 是从 E 到 E 上的双射，由此即得断言。

我们来证明 \(b\)）。若 \(\mu \neq \lambda\) 且 \(x \in \mathrm{Ker}((\lambda - \mu)^{-1}1_{\mathrm{E}} - \mathrm{R}(u, \lambda))\)，则 \(x \in \mathrm{dom}(u)\)，且公式 \(1_{\mathrm{E}} = (\lambda 1_{\mathrm{E}} - u) \circ \mathrm{R}(u, \lambda)\) 蕴含 \(x \in \mathrm{Ker}(\mu 1_{\mathrm{E}} - u)\)。反之，若 \(x \in \mathrm{Ker}(\mu 1_{\mathrm{E}} - u)\) 且 \(\mu \neq \lambda\)，则公式 \(1_{\mathrm{dom}(u)} = \mathrm{R}(u, \lambda) \circ (\lambda 1_{\mathrm{E}} - u)\) 蕴含 \(\mathrm{R}(u, \lambda)(x) = (\lambda - \mu)^{-1}x\)。

命题 16. --- 设 u 是复希尔伯特空间 E 上稠定的闭算子。\(u^{*}\) 的谱是 u 的谱在复共轭下的像，且对 u 的预解集的每个元素 \(\lambda\)，有 \(\mathrm{R}(u,\lambda)^{*}=\mathrm{R}(u^{*},\overline{\lambda})\)。特别地，若 u 自伴，则自同态 \(\mathrm{R}(u,\lambda)\) 对每个 \(\lambda\notin\mathrm{Sp}(u)\) 都是正规的。

设 \(\lambda\in\mathbf{C} - \mathrm{Sp}(u)\) 是 \(u\) 的预解集中的一个元素。我们有 \(u=\lambda1_{\mathrm{E}}-\mathrm{R}(u,\lambda)^{-1}\)，于是

\[
u ^ {*} = \overline {{\lambda}} 1 _ {\mathrm{E}} - (\mathrm{R} (u, \lambda) ^ {- 1}) ^ {*} = \overline {{\lambda}} 1 _ {\mathrm{E}} - (\mathrm{R} (u, \lambda) ^ {*}) ^ {- 1}
\]

（IV，第 236 页及 IV，第 239 页命题 9）。由注 5 推出 \(\overline{\lambda} \in \mathbf{C} - \mathrm{Sp}(u^{*})\) 且 \(\mathrm{R}(u, \lambda)^{*} = \mathrm{R}(u^{*}, \overline{\lambda})\)。因此，\(u^{*}\) 的谱包含于 \(\mathrm{Sp}(u)\) 经复共轭所成的像。把该性质应用于 \(u^{*}\) 即得等式，因为 \(u^{**} = u\)（IV，第 237 页命题 8 的推论）。最后的断言则由第 245 页公式 (9) 得出。

推论。--- 设 u 是复希尔伯特空间 E 上的自伴部分算子。若 E 非零，则 u 的谱非空。

设 \(\operatorname{Sp}(u)\) 为空集。则 u 是单射的，且 \(u^{-1} = -\operatorname{R}(u,0)\) 是 E 的单射自同态，满足 \(\operatorname{Sp}(u^{-1}) \subset \{0\}\)（命题 15，a）），从而 \(\operatorname{Sp}(u^{-1}) = \{0\}\)（I，第 26 页，推论 1）。由于 \(u^{-1}\) 正规（命题 16），这就蕴含 \(u^{-1}\) 为零（I，第 110 页，例 1），此为矛盾。

引理 5. --- 设 E 为复希尔伯特空间，u 为 E 上的闭稠定部分算子。设 \(\lambda \in \mathbf{C}\)。假设存在实数 \(c > 0\)，使得

\[
\| u (x) - \lambda x \| \geqslant c \| x \| \quad \text{对一切 } x \in \operatorname{dom} (u),\tag{10}
\]

\[
\| u ^ {*} (x) - \overline {{\lambda}} x \| \geqslant c \| x \| \quad \text{对一切 } x \in \operatorname{dom} (u ^ {*}),\tag{11}
\]

则 \(\lambda\) 属于 \(u\) 的预解集，且 \(\| \mathrm{R}(u,\lambda)\| \leqslant c^{-1}\)。

假设蕴含 \(u - \lambda 1_{\mathrm{E}}\) 与 \(u^{*} - \overline{\lambda} 1_{\mathrm{E}}\) 都是单射的。设 F 为 \(u - \lambda 1_{\mathrm{E}}\) 的像。空间 F 在 E 中稠密，因为其正交补等于 \(\operatorname{Ker}(u^{*} - \overline{\lambda} 1_{\mathrm{E}})\)（IV，第 236 页，命题 7，c）），而后者为零。

我们来证明空间 F 是闭的。设 \((x_{n})_{n\in\mathbf{N}}\) 是 dom(u) 中的序列，使序列 \((u(x_{n})-\lambda x_{n})_{n\in\mathbf{N}}\) 收敛于 \(y\in F\)。不等式 (10) 蕴含序列 \((x_{n})_{n\in\mathbf{N}}\) 是 E 中的柯西序列。设 \(x\in E\) 为其极限。序列 \((x_{n},u(x_{n}))\) 在 E×E 中收敛于 \((x,y+\lambda x)\)；由于 u 的图像是闭的，我们因此有 \(x\in\text{dom}(u)\) 与 \(u(x)=y+\lambda x\)，这就证明了 \(y\in F\)。

由此断定 F = E。于是部分算子 \(u - \lambda1_{E}\) 是双射，从而 \(\lambda \notin \operatorname{Sp}(u)\)（注 3）。不等式 (10) 进而又蕴含 \(\|R(u, \lambda)\| \leqslant c^{-1}\)。

命题 17. — 设 E 为复希尔伯特空间，u 为 E 上的自伴部分算子。

\begin{enumerate}
\def\labelenumi{\alph{enumi})}
\item
  u 的谱包含于 \(\mathbf{R}\)；
\item
  若 u 是正的，则 u 的谱包含于 \(\mathbf{R}_{+}\)；
\item
  设 E 非零。设 \(\lambda \notin \operatorname{Sp}(u)\)，且 \(\delta > 0\) 为 \(\lambda\) 到 \(u\) 的谱的距离。我们有 \(\| \mathrm{R}(u, \lambda) \| = \delta^{-1}\)。
\end{enumerate}

设 \((a,b)\in\mathbf{R}\times\mathbf{R}\)，\(\lambda=a+ib\)。设 \(x\in\mathrm{dom}(u)\)。由于 u 自伴，我们有 \(\langle x|u(x)\rangle\in\mathbf{R}\)，由此

\[
\begin{array}{r} \| u (x) - \lambda x \| ^ {2} = \| u (x) \| ^ {2} - 2 a \langle x | u (x) \rangle + (a ^ {2} + b ^ {2}) \| x \| ^ {2} \\ = \| u (x) - \overline {{\lambda}} x \| ^ {2}. \end{array}
\]

设 \(b \neq 0\)。于是得到

\[
\begin{array}{c} \| u (x) - \lambda x \| ^ {2} = \| u (x) - \overline {{\lambda}} x \| ^ {2} \geqslant (\| u (x) \| - a \| x \|) ^ {2} + b ^ {2} \| x \| ^ {2} \\ \geqslant b ^ {2} \| x \| ^ {2}. \end{array}
\]

于是由引理 5 得 \(\lambda \notin \operatorname{Sp}(u)\)，即断言 \(a\)）。

设 u 还是正的。若 b = 0 且 a \textless{} 0，则对 \(x \in \text{dom}(u)\) 类似地得到不等式

\[
\left\| u (x) - \lambda x \right\| ^ {2} = \left\| u (x) - \overline {{\lambda}} x \right\| ^ {2} \geqslant \left(\left\| u (x) \right\| - a \| x \|\right) ^ {2} \geqslant a ^ {2} \| x \| ^ {2},
\]

从而 \(\lambda\notin\mathrm{Sp}(u)\)（同上所引），这就证明了 \(b\)）。

最后证明 \(c\)）。由命题 16，预解式 \(\mathrm{R}(u,\lambda)\) 是 E 的正规自同态。因此其范数等于其谱半径（I，第 108 页，推论 1），由此

\[
\| \mathrm{R} (u, \lambda) \| = \sup _ {\mu \in \operatorname{Sp} (\mathrm{R} (u, \lambda))} | \mu |
\]

（I，第 24 页，定理 1）。R(u,λ) 的谱不可能化为 \{0\}，因为否则将有 \(\| \mathrm{R}(u,\lambda)\| = 0\)，于是作为 R(u,λ) 之像的 \(\operatorname{dom}(u)\) 将为零，E 亦然。因此命题 15 蕴含

\[
\| \mathrm{R} (u, \lambda) \| = \sup _ {\mu \in \mathrm{Sp} (u)} \frac {1}{| \lambda - \mu |} = \frac {1}{\delta}.
\]

推论。--- 设 u 是 E 上的自伴算子。

\begin{enumerate}
\def\labelenumi{\alph{enumi})}
\item
  u 的剩余谱是空集；
\item
  对所有 \(\lambda \neq \mu\)（在 \(\mathbf{C}\) 中），\(u\) 的属于 \(\lambda\) 与 \(\mu\) 的特征空间相互正交。
\end{enumerate}

我们来证明 a)。设 \(\lambda\) 属于 \(u\) 的谱；它是实数（命题 17）。我们有 \(\mathrm{Ker}(\lambda 1_{\mathrm{E}} - u) = \mathrm{Im}(\lambda 1_{\mathrm{E}} - u)^{\circ}\)（IV，第 236 页，命题 7），因此 \(\lambda\) 不是 \(u\) 的特征值，只要部分算子 \(\lambda 1_{\mathrm{E}} - u\) 的像稠密。按定义，这就蕴含 \(u\) 的剩余谱是空集。

我们来证明 \(b\)）。由该命题，可以假设 \(\lambda\) 与 \(\mu\) 是实数。设 \(x \in \operatorname{dom}(u)\) 满足 \(u(x) = \lambda x\)，且 \(y \in \operatorname{dom}(u)\) 满足 \(u(y) = \mu y\)。于是有

\[
\lambda \langle x \mid y \rangle = \langle u (x) \mid y \rangle = \langle x \mid u (y) \rangle = \mu \langle x \mid y \rangle ,
\]

由此 \(\langle x \mid y \rangle = 0.\)

注。--- 若 u 是非自伴的闭对称算子，则其谱不包含于 R（参见下文 IV，第 257 页推论 10），而且 u 的对应于 λ 与 \(\overline{\lambda}\) 的特征空间有可能不正交（IV，第 347 页，习题 11）。

\subsection{伪谱}

定义 9. --- 设 E 为复巴拿赫空间，u 为 E 上的闭部分算子。设 \(\varepsilon\) 为严格正的实数。我们把 u 的 \(\varepsilon\)-伪谱称为 u 的谱与如下的复数 \(\lambda\) 的集合之并：这些复数属于 u 的预解集且满足 \(\| \mathrm{R}(u,\lambda)\| >\varepsilon^{-1}\)。我们把这个集合记作 \(\mathrm{PSp}_{\varepsilon}(u)\)。

\pandocbounded{\includegraphics[keepaspectratio]{images/part002_14758b8b.jpg}}

有些作者把 \(\varepsilon\)-伪谱（对 \(u\)）定义为集合 \(\mathrm{T}_{\varepsilon}(u)\)，即 \(\mathrm{Sp}(u)\) 与如下的 \(\lambda \in \mathbf{C} - \mathrm{Sp}(u)\) 所成集合之并：它们满足 \(\| \mathrm{R}(u,\lambda)\| \geqslant \varepsilon^{-1}\)。\(\mathrm{PSp}_{\varepsilon}(u)\) 的闭包包含于 \(\mathrm{T}_{\varepsilon}(u)\)，但这一包含可以是严格的，即使 E 是希尔伯特空间亦然（参见 IV，第 348 页习题 18 及第 349 页习题 19）。

命题 18. --- 设 E 为复巴拿赫空间，u 为 E 上的闭部分算子。设 \(\varepsilon \in \mathbf{R}_{+}^{*}\)。集合 \(\mathrm{PSp}_{\varepsilon}(u)\) 是 C 的开子集。当 E 非零时它非空。

若 E 为零，则 \(\mathrm{PSp}_{\varepsilon}(u)\) 是空集。设 E 非零。由 IV，第 246 页命题 14，集合 \(\mathrm{PSp}_{\varepsilon}(u)\) 是 C 的开子集。

集合 \(\mathrm{PSp}_{\varepsilon}(u)\) 非空，如果 \(u\) 的谱非空。若 \(\mathrm{Sp}(u)\) 为空集，则刘维尔定理（VAR，R1，第 29 页，3.3.6）蕴含 \(\mathbf{C}\) 上由 \(\lambda \mapsto \mathrm{R}(u,\lambda)\) 定义的全纯函数不有界，故存在 \(\lambda\) 在 \(\mathbf{C}\) 中，使得 \(\| \mathrm{R}(u,\lambda)\| >\varepsilon^{-1}\)。

命题 19. --- 设 E 为复巴拿赫空间，u 为 E 上的闭部分算子。则我们有 \(\mathrm{PSp}_{\varepsilon}(u) \subset \mathrm{PSp}_{\delta}(u)\)（当 \(0 < \varepsilon < \delta\) 时），并且

\[
\bigcap_ {\varepsilon \in \mathbf {R} _ {+} ^ {*}} \mathrm{PSp} _ {\varepsilon} (u) = \mathrm{Sp} (u).
\]

第一个断言由定义立得。至于第二个断言，按定义，u 的谱包含于 \(\mathrm{PSp}_{\varepsilon}(u)\)，对一切 \(\varepsilon > 0\) 成立；而若 \(\lambda \notin \mathrm{Sp}(u)\)，则 \(\lambda \notin \mathrm{PSp}_{\varepsilon}(u)\)，只要 \(\varepsilon < \|R(u, \lambda)\|^{-1}\)。

命题 20. --- 设 E 为复巴拿赫空间，u 为 E 上的闭部分算子。设 \(\varepsilon\in\mathbf{R}_{+}^{*}\)。对每个 \(\lambda\in\mathbf{C}\)，下列条件等价：

\begin{enumerate}
\def\labelenumi{(\roman{enumi})}
\item
  有 \(\lambda \in \mathrm{PSp}_{\varepsilon}(u)\)；
\item
  或者 \(\lambda \in \mathrm{Sp}(u)\)，或者存在 \(x \in \mathrm{dom}(u)\) 满足 \(\| x \| = 1\) 且 \(\| (\lambda 1_{\mathrm{E}} - u)(x) \| < \varepsilon\)；
\item
  存在 \(v \in \mathcal{L}(\mathrm{E})\) 满足 \(\| v \| < \varepsilon\) 且 \(\lambda \in \mathrm{Sp}(u + v)\)。
\end{enumerate}

可以假设 E 非零。条件 (ii) 是定义的重新表述，因而也就是条件 (i)。

设条件 (i) 成立，我们来证明 (iii)。若 \(\lambda\) 属于 \(u\) 的谱，则在 (iii) 中可取 \(v = 0\)。

于是设 \(\lambda \notin \mathrm{Sp}(u)\)。由 \(\mathrm{PSp}_{\varepsilon}(u)\) 的定义，存在 \(y \in \mathrm{E}\) 满足 \(\| y \| = 1\) 且 \(\| \mathrm{R}(u, \lambda)y \| > \varepsilon^{-1}\)。令 \(x = \mathrm{R}(u, \lambda)y\)。我们有 \(x \neq 0\)。由哈恩–巴拿赫定理（EVT，II，第 67 页，推论 2），线性形式 \(\ell\)（在 \(\mathbf{C}x\) 上、满足 \(\ell(x) = 1\)）容许连续延拓 \(\ell_1 \in \mathrm{E}'\) 满足 \(\| \ell_1 \| = \| \ell \|\)；于是 \(\| \ell_1 \| = \| x \|^{-1} < \varepsilon\)。对每个 \(e \in \mathrm{E}\)，令 \(v(e) = \ell_1(e)y\)。我们有 \(v \in \mathcal{L}(\mathrm{E})\) 且 \(v(x) = y\)。我们得到 \((u + v)x = u(x) + y = \lambda x\)，故 \(\lambda \in \mathrm{Sp}(u + v)\)。又因 \(\| v \| = \| \ell_1 \| < \varepsilon\)，条件 (iii) 得到满足。

反之，设 \(v \in \mathcal{L}(\mathrm{E})\) 满足 \(\|v\| < \varepsilon\) 且 \(\lambda \in \mathrm{Sp}(u + v)\)。于是部分算子 \(\lambda1_{\mathrm{E}} - (u + v)\) 或者不是单射的，或者没有连续的逆；把 IV，第 243 页的引理 4 应用于单射部分算子 \(\lambda1_{E} - u\) 与 \(-v\)，便得 \(\|v\| \geqslant \|R(u, \lambda)\|^{-1}\)。由此推出条件 (iii) 蕴含 (i)。

推论。— 设 E 为复巴拿赫空间，u 为 E 上的闭部分算子。设 \(\varepsilon > 0\)。

\begin{enumerate}
\def\labelenumi{\alph{enumi})}
\item
  对每个 \(v \in \mathcal{L}(\mathrm{E})\)，我们有 \(\mathrm{PSp}_{\varepsilon}(u) \subset \mathrm{PSp}_{\varepsilon + \| v\|}(u + v)\)；
\item
  设 \(\delta > 0\)，\(D_{\delta}\) 为以 0 为心、以 \(\delta\) 为半径的 \(\mathbf{C}\) 中的开圆盘。我们有 \(\mathrm{PSp}_{\varepsilon}(u) + D_{\delta} \subset \mathrm{PSp}_{\varepsilon + \delta}(u)\)。
\end{enumerate}

设 \(\lambda \in \mathrm{PSp}_{\varepsilon}(u)\)。存在自同态 \(w \in \mathcal{L}(\mathrm{E})\) 满足 \(\|w\| < \varepsilon\) 且 \(\lambda \in \mathrm{Sp}(u+w)\)（命题 20）。由于 \(u+w = (u+v)+(w-v)\) 且 \(\|w-v\| < \varepsilon + \|v\|\)，我们有 \(\lambda \in \mathrm{PSp}_{\varepsilon+ \|v\|}(u+v)\)（同上所引）。

设 \(\mu \in \mathrm{D}_{\delta}\)；我们有 \(\lambda + \mu \in \mathrm{Sp}(u + (w + \mu 1_{\mathrm{E}}))\) 且 \(\| w + \mu 1_{\mathrm{E}} \| < \varepsilon + \delta\)，故 \(\lambda + \mu \in \mathrm{PSp}_{\varepsilon + \delta}(u)\)（同上所引）。

命题 21. --- 设 E 为复巴拿赫空间，u 为 E 上的闭部分算子。设 \(\varepsilon\in\mathbf{R}_{+}^{*}\)。\(\mathrm{PSp}_{\varepsilon}(u)\) 的每个有界连通分支都与 \(u\) 的谱相交。

设 U 是 \(\mathrm{PSp}_{\varepsilon}(u)\) 的一个不与 \(\mathrm{Sp}(u)\) 相交的连通分支。集合 U 在 \(\mathrm{PSp}_{\varepsilon}(u)\) 中既开又闭，因此它在 C 中的闭包 \(\overline{U}\) 满足 \(\overline{U} \cap \mathrm{PSp}_{\varepsilon}(u) = U\)。

由于 \(\mathrm{Sp}(u)\) 包含于 \(\mathrm{PSp}_{\varepsilon}(u)\) 而 U 与 \(\mathrm{Sp}(u)\) 不相交，这一等式表明集合 \(\overline{U}\) 与 \(\mathrm{Sp}(u)\) 不相交，从而包含于 u 的预解集。此外，它还蕴含集合 \(\overline{U}-U\) 与 \(\mathrm{PSp}_{\varepsilon}(u)\) 不相交。

于是我们有 \(\|\mathrm{R}(u,\lambda)\| \leqslant \varepsilon^{-1}\)，对一切 \(\lambda\)（属于 \(\overline{U} - U\)）；而 \(\|\mathrm{R}(u,\lambda)\| > \varepsilon^{-1}\)，对 \(\lambda \in U\)。若集合 U 有界，则其闭包 \(\overline{U}\) 是紧的，且存在 \(\lambda_{0} \in \overline{U}\) 使得 \(\|\mathrm{R}(u,\lambda)\| \leqslant \|\mathrm{R}(u,\lambda_{0})\|\) 对 \(\lambda \in \overline{U}\) 成立。前述结果表明 \(\lambda_{0} \in U\)，这与极大值原理（VAR，R1，第 29 页，3.3.7）矛盾，因为 u 的预解式在 u 的预解集中是全纯的（IV，第 246 页，命题 14）。

\subsection{乘法算子}

设 X 是局部紧空间，\(\mu\) 是 X 上的正测度。我们考察 \(\mathrm{L}^{2}(\mathrm{X},\mu)\) 上的乘法算子；它们是稠定的闭算子（IV，第 232 页，命题 5）。对 X 上任何 \(\mu\)-可测函数 g，我们都用 \(m_{g}\) 表示 \(\mathrm{L}^{2}(\mathrm{X},\mu)\) 上以 g 为乘子的乘法部分算子。

命题 22. — 设 \(g\) 是 X 上的 \(\mu\)-可测函数。

\begin{enumerate}
\def\labelenumi{\alph{enumi})}
\item
  \(m_{g}\) 的谱是相对于 \(\mu\) 的 \(g\) 的本质像 S；
\item
  设 \(\lambda \in \mathbf{C} - \mathrm{Sp}(m_g)\)。预解式 \(\mathrm{R}(m_g, \lambda)\) 是乘法算子 \(m_h\)，其中 \(h\) 是 X 上如下定义的函数：\(h(x) = 0\)（若 \(g(x) = \lambda\)），否则 \(h(x) = (\lambda - g(x))^{-1}\)。
\end{enumerate}

我们来证明 \(\mathbf{C} - \mathbf{S}\) 包含于 \(m_g\) 的预解集。设 \(\lambda \in \mathbf{C} - \mathbf{S}\)。存在 \(\lambda\) 的开邻域 U，使 X 的子集 \(\mathrm{Y} = g^{-1}(\mathrm{U})\) 是局部 \(\mu\)-零测的。X 上如下定义的函数 \(k\)：\(k(x) = (\lambda - g(x))^{-1}\)（若 \(x \notin \mathrm{Y}\)），而 \(k(x) = 0\)（若 \(x \in \mathrm{Y}\)），它属于 \(\mathcal{L}^{\infty}(\mathrm{X}, \mu)\)（IV，第 184 页，引理 5）；因此以 \(k\) 为乘子的乘子算子是 \(\mathrm{L}^{2}(\mathrm{X}, \mu)\) 的自同态。

由于 \(|gk| \leqslant 1 + |\lambda k|\)，我们有

\[
| g k f | \leqslant | f | + | \lambda k f |
\]

对 \(f \in \mathcal{L}^2(\mathrm{X},\mu)\) 成立，这蕴含 \(m_k\) 的像包含于 \(m_g\) 的定义域。反之，设 \(f \in \mathcal{L}^2(\mathrm{X},\mu)\)，其类 \(\widetilde{f}\) 属于 \(m_g\) 的定义域。在局部 \(\mu\)-零测集 Y 之外，我们有 \(f(x) = k(x)(\lambda - g(x))f(x)\)，故 \(\widetilde{f}\) 在 \(m_k\) 的像中。同一公式证明 \(\lambda\) 属于 \(m_g\) 的预解集且 \(\mathrm{R}(m_g,\lambda) = m_k\)。由于集合 Y 是局部 \(\mu\)-零测的，乘法算子 \(m_k\) 与算子 \(m_h\)（如断言 \(b\)）中所描述的）一致。

反之，我们来证明 \(\mathbf{C} - \mathrm{Sp}(m_g)\) 包含于 \(\mathbf{C} - \mathrm{S}\)。设 \(\lambda \in \mathbf{C} - \mathrm{Sp}(m_g)\)。取实数 \(\mathrm{M} > \| \mathrm{R}(m_g, \lambda) \|\)。用 Y 表示由这样的 \(x \in \mathrm{X}\) 组成的集合：它们满足 \(|\lambda - g(x)| < \mathrm{M}^{-1}\)。我们来证明 Y 是局部 \(\mu\)-零测的，这将蕴含 \(\lambda\) 不属于 S，并完成证明。

设 \(K\) 是 \(X\) 的紧子集。设 \(\varphi\) 是 \(Y \cap K\) 的特征函数；它是 \(\mathcal{L}^{2}(X,\mu)\) 中的元素，我们把它在 \(\mathrm{L}^{2}(X,\mu)\) 中的类记作 \(\widetilde{\varphi}\)。设 \(\psi\) 是 \(\mathcal{L}^{2}(X,\mu)\) 中的函数，它在 \(\mathrm{L}^{2}(X,\mu)\) 中的类是 \(\mathrm{R}(m_g,\lambda)(\widetilde{\varphi})\)。我们有 \(\mathrm{R}(m_g,\lambda)(\widetilde{\varphi}) \in \operatorname{dom}(m_g)\) 且 \((\lambda - m_g)(\mathrm{R}(m_g,\lambda)(\widetilde{\varphi})) = \widetilde{\varphi}\)，故 \((\lambda - g(x))\psi(x) = 1\) 对 \(\mu\)-几乎所有的 \(x \in Y \cap K\) 成立。由此可得

\[
\| \mathrm{R} (m _ {g}, \lambda) \| ^ {2} \| \widetilde {\varphi} \| ^ {2} \geqslant \| \mathrm{R} (m _ {g}, \lambda) (\widetilde {\varphi}) \| ^ {2} \geqslant \mathrm{M} ^ {2} \mu (\mathrm{Y} \cap \mathrm{K}) = \mathrm{M} ^ {2} \| \widetilde {\varphi} \| ^ {2}.
\]

鉴于 M 的取法，这表明 \(\varphi\) \(\mu\)-几乎处处为零。于是集合 Y \(\cap\) K 是 \(\mu\)-零测的。因此集合 Y 是局部 \(\mu\)-零测的（INT，IV，第 172 页，§ 5，第 2 目，命题 5）。

命题 23. --- 设 g 是 X 上的 μ-可测函数。乘法算子 \(m_{g}\) 的伴随算子是 \(m_{\overline{g}}\)。

对每个整数 \(n \geqslant 1\)，设 \(\varphi_{n}\) 是由这样的元素 \(x \in X\) 组成的集合的特征函数：它们满足 \(|g(x)| \leqslant n\)，并设 \(\widetilde{\varphi}_{n}\) 是它在 \(\mathrm{L}^{2}(\mathrm{X}, \mu)\) 中的类。设 \(f \in \mathcal{L}^{2}(\mathrm{X}, \mu)\)，其类 \(\widetilde{f}\) 属于 \(\operatorname{dom}(m_{g}^{*})\)，又设 \(\psi\) 是以 \(m_{g}^{*}(\widetilde{f})\) 为类的函数。

对每个 \(h \in \mathcal{L}^2(\mathrm{X}, \mu)\)，若其类 \(\widetilde{h}\) 属于 \(\operatorname{dom}(m_g)\)，则也有 \(\widetilde{h}\widetilde{\varphi}_n \in \operatorname{dom}(m_g)\)，从而 \(\langle \widetilde{f} | m_g(\widetilde{h}\widetilde{\varphi}_n) \rangle = \langle m_g^*(\widetilde{f}) | \widetilde{h}\widetilde{\varphi}_n \rangle\)。由此得到等式

\[
\int_ {\mathrm{X}} \overline {{(f \overline {{g}} - \psi)}} \varphi_ {n} h d \mu = 0.
\]

由于 \(m_{g}\) 的定义域在 \(\mathrm{L}^{2}(\mathrm{X},\mu)\) 中稠密，我们推出 \((f\overline{g}-\psi)\varphi_{n}\) \(\mu\)-几乎处处为零。由于 n 是任意的，这就意味着 \(m_{g}^{*}(\widetilde{f})\) 在 \(\mathrm{L}^{2}(\mathrm{X},\mu)\) 中等于 \(f\overline{g}\) 的类。特别地，由 \(m_{g}^{*}(\widetilde{f}) \in \mathrm{L}^{2}(\mathrm{X},\mu)\)，我们断定 f 属于 \(m_{\overline{g}}\) 的定义域，且 \(m_{g}^{*}(\widetilde{f}) = m_{\overline{g}}(\widetilde{f})\)。

因此 \(m_{g}\) 的伴随算子是 \(m_{\overline{g}}\) 的扩张。此外，我们有

\[
\langle f \mid m _ {g} (h) \rangle = \int_ {\mathrm{X}} \overline {{f}} \cdot (g h) d \mu = \int_ {\mathrm{X}} f \overline {{g}} h d \mu
\]

对一切 \(f \in \mathrm{L}^{2}(\mathrm{X}, \mu)\) 与 \(h \in \mathrm{dom}(m_{g})\) 成立，这就证明线性形式 \(h \mapsto \langle f | m_{g}(h) \rangle\) 当 \(m_{\overline{g}}(f) g\) 属于 \(\mathrm{L}^{2}(\mathrm{X}, \mu)\) 时是连续的。于是，\(m_{\overline{g}}\) 的定义域包含于 \(m_{g}^{*}\) 的定义域，证明至此完成。

推论。--- 设 g 是 X 上的 μ-可测函数。乘法算子 \(m_{g}\)（\(\mathrm{L}^{2}(\mathrm{X},\mu)\) 上的）是自伴的（相应地，正的），当且仅当函数 g 局部 μ-几乎处处取实值（相应地，局部 μ-几乎处处为正）。

第一个断言由该命题得出。若 \(g\) 局部 \(\mu\)-几乎处处为正，则有 \(\langle f|m_{g}(f)\rangle=\int_{\mathrm{X}}g|f|^{2}d\mu\geqslant0\)，对一切 \(f\in\mathrm{L}^{2}(\mathrm{X},\mu)\) 成立，故部分算子 \(m_{g}\) 是正的。

反之，若 \(m_{g}\) 是正的，则其谱包含于 \(\mathbf{R}_{+}\)（IV，第 248 页，命题 17）；由于该谱是 g 的 \(\mu\)-本质像（命题 22，a）），这就意味着 g 局部 \(\mu\)-几乎处处为正。

引理 6. --- 设 \(g_{1}\) 与 \(g_{2}\) 是 X 上的 \(\mu\)-可测函数。

\begin{enumerate}
\def\labelenumi{\alph{enumi})}
\item
  部分算子 \(m_{g_{1}} + m_{g_{2}}\) 是可闭的，其闭包是乘法算子 \(m_{g_{1}+g_{2}}\)；
\item
  有 \(m_{g_{1}} \circ m_{g_{2}} \subset m_{g_{1}g_{2}}\)；
\item
  设 \(g_{1}\) 有界。这时我们有 \(m_{g_{2}} \circ m_{g_{1}} = m_{g_{1}g_{2}}\)。此外，\(m_{g_{2}}\) 的定义域包含于 \(m_{g_{1}g_{2}}\) 的定义域，且 \(m_{g_{1}} \circ m_{g_{2}}\) 是 \(m_{g_{1}g_{2}}\) 到 \(\text{dom}(m_{g_{2}})\) 上的限制。
\end{enumerate}

\(m_{g_{1}+g_{2}}\) 是 \(m_{g_{1}} + m_{g_{2}}\) 的扩张，这是初等的；于是后一算子是可闭的，且 \(\overline{m_{g_{1}} + m_{g_{2}}} \subset m_{g_{1}+g_{2}}\)。

设 \(f \in \mathcal{L}^2(\mathrm{X},\mu)\)，使函数 \(h = (g_1 + g_2)f\) 属于 \(\mathcal{L}^2(\mathrm{X},\mu)\)。对每个整数 \(n \geqslant 1\)，用 \(\mathrm{X}_n\) 表示由这样的 \(x \in \mathrm{X}\) 组成的集合：它们满足 \(|g_1(x)| + |g_2(x)| \leqslant n\)，并设 \(\varphi_n\) 是 \(\mathrm{X}_n\) 的特征函数。我们有 \(\varphi_n f \in \mathrm{dom}(m_{g_1} + m_{g_2})\)。由于 \((\varphi_n f)(x)\) 收敛于 \(f(x)\)，对一切 \(x \in \mathrm{X}\)，且对每个 \(n \in \mathbf{N}\) 有 \(|\varphi_n f| \leqslant |f|\) 与 \(|(g_1 + g_2)\varphi_n f| \leqslant |(g_1 + g_2)f| = |h|\)（由假设 \(h \in \mathcal{L}^2(\mathrm{X},\mu)\)），勒贝格定理（INT，IV，第 137 页，§ 3，第 7 目，定理 6）蕴含：\(\mathrm{L}^{2}(\mathrm{X},\mu)\times\mathrm{L}^{2}(\mathrm{X},\mu)\) 中 \((\varphi_{n}f,(g_{1}+g_{2})\varphi_{n}f)\) 的类对（它们属于 \(m_{g_{1}}+m_{g_{2}}\) 的图像）所成的序列，收敛于 \((f,h)\) 在 \(\mathrm{L}^{2}(\mathrm{X},\mu)\) 中的类对。因此 \(m_{g_{1}}+m_{g_{2}}\) 的闭包确实等于 \(m_{g_{1}+g_{2}}\)。

\(m_{g_{1}} \circ m_{g_{2}} \subset m_{g_{1}g_{2}}\) 是初等的。设 \(g_{1}\) 有界，于是 \(m_{g_{1}}\) 是 \(\mathrm{L}^{2}(\mathrm{X},\mu)\) 上的连续线性映射。此时 \(m_{g_{2}} \circ m_{g_{1}}\) 的定义域是这样的函数 \(f \in \mathcal{L}^{2}(\mathrm{X},\mu)\) 的等价类所成的集合：这些函数使 \(g_{2}(g_{1}f) \in \mathcal{L}^{2}(\mathrm{X},\mu)\)，亦即 \(m_{g_{1}g_{2}}\) 的定义域。因此我们有 \(m_{g_{2}} \circ m_{g_{1}} = m_{g_{1}g_{2}}\)。

同样初等的是：\(\operatorname{dom}(m_{g_{1}} \circ m_{g_{2}}) = \operatorname{dom}(m_{g_{2}})\) 包含于 \(\operatorname{dom}(m_{g_{1}g_{2}})\)，且 \(m_{g_{1}g_{2}}\) 到这个空间的限制等于 \(m_{g_{1}} \circ m_{g_{2}}\)。

应当警惕，不要以为一般地有 \(m_{g_{1}} \circ m_{g_{2}} = m_{g_{1}g_{2}}\)（IV，第 347 页，习题 10）。尽管如此，我们有如下的部分结果：

命题 24. --- 设 g 是 X 上的 μ-可测函数。

\begin{enumerate}
\def\labelenumi{\alph{enumi})}
\item
  有 \(m_{\overline{g}} \circ m_g = m_{|g|^2}\)；
\item
  对一切整数 \(k, \ell \in \mathbf{N}\)，有 \(m_{g^k\overline{g}^\ell} = m_g^k m_{\overline{g}}^\ell\)。
\end{enumerate}

由引理 6 有 \(m_{\overline{g}} \circ m_g \subset m_{|g|^2}\)。反之，由不等式 \(|g| \leqslant 1 + |g|^2\) 推出 \(\text{dom}(m_{|g|^2}) = \text{dom}(m_{\overline{g}} \circ m_g)\)，由此得第一个断言。

设 \(k, \ell \in \mathbf{N}\)。我们有 \(m_g^k m_{\overline{g}}^\ell \subset m_{g^k \overline{g}^\ell}\)（同上所引）。\(m_g^k m_{\overline{g}}^\ell\) 的定义域是在 \(\mathrm{L}^2(\mathrm{X}, \mu)\) 中这样的函数 \(h \in \mathcal{L}^2(\mathrm{X}, \mu)\) 的类所成的集合：\(|g|^j h\) 属于 \(\mathcal{L}^2(\mathrm{X}, \mu)\)，对每个整数 \(j\)（满足 \(0 \leqslant j \leqslant k + \ell\)）。不等式

\[
\left| g ^ {j} h \right| \leqslant | h | + \left| g ^ {k + \ell} h \right|,
\]

对 \(0 \leqslant j \leqslant k + \ell\) 成立，由它可以看出 \(\mathrm{dom}(m_g^k m_{\overline{g}}^\ell)\) 等于 \(\mathrm{dom}(m_{g^k\overline{g}^\ell})\)，由此得断言 \(b\)）。

\subsection{对称算子的自伴扩张}

在本节中，我们打算对复希尔伯特空间上对称算子的自伴扩张进行分类，特别是找出保证这类算子的自伴扩张存在的条件。

设 E 是复希尔伯特空间。对 E 上每个部分算子 u，本节中我们将把 \(u + i\) 与 \(u - i\) 分别写作 \(u + i1_{E}\) 与 \(u - i1_{E}\) 的简写。

设 u 是 E 上稠定的闭算子。我们回顾（参见 IV，第 230 页定义 4 及 IV，第 230 页命题 4）：我们用 \(E_{u}\) 表示赋有埃尔米特形式的希尔伯特空间 dom(u)

\[
(x \mid y) _ {u} = \langle x \mid y \rangle + \langle u (x) \mid u (y) \rangle .
\]

我们用 \(\|x\|_{u}\) 表示希尔伯特空间 \(E_{u}\) 中元素 x 的范数。dom(u) 的每个在 E 中闭的子空间都在 \(E_{u}\) 中闭（IV，第 230 页，注 1）。

定义 10. --- 设 u 是复希尔伯特空间 E 上的闭对称算子。设 \(\mathrm{E}_{+} = \mathrm{Ker}(u^{*} - i)\)，\(\mathrm{E}_{-} = \mathrm{Ker}(u^{*} + i)\)。子空间偶 \((\mathrm{E}_{+}, \mathrm{E}_{-})\)（\(\mathrm{dom}(u^{*})\) 的）称为 u 的亏偶。子空间 \(E_{+}\) 与 \(E_{-}\) 是 E 的闭子空间（IV，第 229 页，命题 3）。这些子空间的希尔伯特维数所成的偶 \((\dim(\mathrm{E}_{+}), \dim(\mathrm{E}_{-}))\) 称为 u 的亏指标。

命题 25. --- 设 u 是复希尔伯特空间 E 上的闭对称算子。我们有

\[
\begin{array}{l} \operatorname{Ker} (u ^ {*} - i) ^ {\circ} = \operatorname{Im} (u + i), \\ \operatorname{Ker} (u ^ {*} + i) ^ {\circ} = \operatorname{Im} (u - i). \end{array}
\]

由 IV，第 236 页命题 7，c)，只需证明 \(u + i\)（相应地，\(u - i\)）的像在 E 中闭。由于 \(u\) 对称，我们有 \(\langle (u + i)(x) \mid (u + i)(y) \rangle = (x \mid y)_u\)，对一切 \(x\) 与 \(y\)（\(\text{dom}(u)\) 中的）。于是映射 \(x \mapsto u(x) + ix\)（从 \(E_u\) 到 E 中的）是等距的。它的像在 E 中闭（I，第 107 页，引理 8）。由于这也是 \(u + i\) 的像，故后者在 E 中闭。同理，\(u - i\) 的像在 E 中闭。

引理 7. --- 设 u 是复希尔伯特空间 E 上的闭对称算子，v 是 u 的闭对称扩张。u 的定义域是希尔伯特空间 \(E_{v}\) 的闭子空间。

由于 v 是 u 的扩张，我们有 \((x \mid y)_{v} = (x \mid y)_{u}\)，对 \(\text{dom}(u)\) 中的 x 与 y。于是 \(E_{u}\) 到 \(E_{v}\) 中的典范单射是等距的，结论由此得出（I，第 107 页，引理 8）。

命题 26. --- 设 u 是复希尔伯特空间 E 上的闭对称算子。设 \((\mathrm{E}_{+}, \mathrm{E}_{-})\) 是 u 的亏偶。空间 \(E_{+}\)、\(E_{-}\) 与 dom(u) 是 \(E_{u^{*}}\) 中两两正交的闭子空间，其希尔伯特直和等于 \(E_{u^{*}}\)。

由于 u 对称，我们有 \(u \subset u^{*}\)，故空间 \(\text{dom}(u)\) 在 \(E_{u^{*}}\) 中闭（引理 7）。子空间 \(E_{+}\) 与 \(E_{-}\) 在 E 中闭且包含于 \(\text{dom}(u^{*})\)，故在 \(E_{u^{*}}\) 中闭。

设 \(x \in \mathrm{E}_{+}\)。我们有 \(u^{*}(x) = ix\)，故 \(\langle x | u(y) \rangle = -i\langle x | y \rangle\) 对一切 \(y \in \operatorname{dom}(u)\) 成立。于是，对每个 \(y \in \operatorname{dom}(u)\)，有

\[
(x \mid y) _ {u ^ {*}} = \langle x \mid y \rangle + \langle u ^ {*} (x) \mid u ^ {*} (y) \rangle = i (\langle x \mid u (y) \rangle - \langle x \mid u ^ {*} (y) \rangle) = 0
\]

因为 u 对称。因此空间 \(E_{+}\) 与 \(\operatorname{dom}(u)\) 在 \(E_{u^{*}}\) 中正交。同理，\(E_{-}\) 与 \(\operatorname{dom}(u)\) 在 \(E_{u^{*}}\) 中正交。

设 \(x \in \mathrm{E}_{+}\)，\(y \in \mathrm{E}_{-}\)。则

\[
(x \mid y) _ {u ^ {*}} = \langle x \mid y \rangle + \langle u ^ {*} (x) \mid u ^ {*} (y) \rangle = \langle x \mid y \rangle - \langle x \mid y \rangle = 0,
\]

故 \(\mathrm{E}_{+}\) 与 \(\mathrm{E}_{-}\) 在 \(\mathrm{E}_{u^{*}}\) 中正交。

设 \(x \in \mathrm{E}_{u^*}\) 与闭子空间 \(\operatorname{dom}(u) \oplus \mathrm{E}_+ \oplus \mathrm{E}_-\) 正交。对每个 \(y \in \operatorname{dom}(u)\)，有 \(u^*(y) = u(y)\)（因为 \(u \subset u^*\)），于是

\[
0 = (x \mid y) _ {u ^ {*}} = \langle x \mid y \rangle + \langle u ^ {*} (x) \mid u ^ {*} (y) \rangle = \langle x \mid y \rangle + \langle u ^ {*} (x) \mid u (y) \rangle .
\]

这蕴含 \(z = u^{*}(x)\) 属于 \(u^{*}\) 的定义域且满足 \(u^{*}(z) = -x\)。令 \(x_{-} = z - ix\)。它是 \(\text{dom}(u^{*})\) 中满足 \(u^{*}(x_{-}) = -x - iz = -ix_{-}\) 的元素。故 \(x_{-}\) 属于 \(E_{-}\)。但是，对每个 \(w \in E_{-}\)，有

\[
\begin{array}{c} - i \langle x _ {-} | w \rangle = - i \langle z | w \rangle + \langle x | w \rangle \\ = \langle u ^ {*} (x) | u ^ {*} (w) \rangle + \langle x | w \rangle = (x | w) _ {u ^ {*}} = 0. \end{array}
\]

取 \(w = x_{-}\)，便得 \(x_{-} = 0\)，即 \(z = ix\)；由于 \(z = u^{*}(x)\)，我们有 \(x \in \mathrm{E}_{+}\)，从而 \(x = 0\)，因为 \(x\) 与 \(\mathrm{E}_{+}\) 在 \(\mathrm{E}_{u^{*}}\) 中正交。这就完成了直和 \(\operatorname{dom}(u) \oplus \mathrm{E}_{+} \oplus \mathrm{E}_{-}\) 等于 \(\mathrm{E}_{u^{*}}\) 的证明。

推论。--- 设 u 是 E 上的闭对称算子。则 u 自伴，当且仅当 u 的亏指标等于 (0,0)。

事实上，对称算子 \(u\) 自伴当且仅当 \(\text{dom}(u^*) = \text{dom}(u)\)；而命题 26 表明 \(\text{dom}(u^*)\) 是 \(\text{dom}(u)\) 与 \(u\) 的亏子空间的希尔伯特直和。

\pandocbounded{\includegraphics[keepaspectratio]{images/part002_f39b57bf.jpg}}

设 \(u\) 是闭对称部分算子。\(u\) 的谱包含于 \(\mathbf{R}\)，当且仅当 \(u\) 自伴。事实上，我们知道 \(\mathrm{Sp}(u) \subset \mathbf{R}\) 当 \(u\) 自伴时成立（IV，第 248 页，命题 17）；若 \(u\) 不自伴，则子空间 \(\mathrm{Ker}(u^{*} + i)\) 与 \(\mathrm{Ker}(u^{*} - i)\) 之一非零，于是 \(u + i\) 或 \(u - i\) 的像是 E 的真子空间（IV，第 256 页，命题 25），从而 \(i\) 或 \(-i\) 属于 \(\mathrm{Sp}(u)\)。

定义 11. --- 设 u 是复希尔伯特空间 E 上的闭对称算子。u 的边界条件是从 \(\operatorname{Ker}(u^{*} - i)\) 到 \(\operatorname{Ker}(u^{*} + i)\) 中的一个部分等距线性映射（EVT，V，第 41 页，定义 3）。

设 \(u\) 是复希尔伯特空间 E 上的闭对称算子，\(b\) 是 \(u\) 的一个边界条件。设 \(I = \text{Ker}(b)^{\circ}\) 为 \(b\) 的初始子空间。我们用 \(u_b\) 表示 \(u^*\) 到 \(\text{dom}(u^*)\) 的如下子空间上的缩减：该子空间是 \(\text{dom}(u)\) 与一个图像的直和，该图像位于 \(E_+ \oplus E_-\) 中，是 \(b\) 在 I 上的限制的图像。由于 \(E_{u^*} = \text{dom}(u) \oplus E_+ \oplus E_-\)（命题 26），部分算子 \(u_b\) 是完全确定的。

换句话说，\(u_{b}\) 的定义域是这样的元素 \(x \in E\) 所成的空间：它们形如 \(x = x_{0} + y + b(y)\)，其中 \(x_{0} \in \text{dom}(u)\)，y 属于 b 的初始子空间。此时我们有 \(u_{b}(x_{0} + y + b(y)) = u(x_{0}) + iy - ib(y)\)。

定理 1（冯·诺伊曼）。--- 设 u 是复希尔伯特空间 E 上的闭对称算子，\((\mathrm{E}_{+}, \mathrm{E}_{-})\) 是它的亏偶。

\begin{enumerate}
\def\labelenumi{\alph{enumi})}
\item
  对 u 的每个边界条件 b，部分算子 \(u_{b}\) 是 u 的闭对称扩张；
\item
  映射 \(b \mapsto u_{b}\) 是从 u 的边界条件集合到 u 的闭对称扩张集合上的双射；
\item
  对 u 的每个边界条件 b，我们有
\end{enumerate}

\[
\begin{array}{l} \operatorname{Im} (u _ {b} + i) = \operatorname{Im} (u + i) \oplus \operatorname{Im} (b) \\ \operatorname{Im} (u _ {b} - i) = \operatorname{Im} (u - i) \oplus \operatorname{Ker} (b) ^ {\circ}; \end{array}
\]

\begin{enumerate}
\def\labelenumi{\alph{enumi})}
\setcounter{enumi}{3}
\tightlist
\item
  对 u 的每个边界条件 b，\(u_{b}\) 的亏偶是 \((\operatorname{Ker}(b), \operatorname{Ker}(b^{*}))\)。
\end{enumerate}

设 \(b\) 是 \(u\) 的一个边界条件，\(\mathrm{I} \subset \mathrm{E}_{+}\) 是它的初始子空间。\(u^{*}\) 到 \(\Gamma_{b|I} \subset \mathrm{E}_{+} \oplus \mathrm{E}_{-}\) 的限制是由 \(x + b(x) \mapsto ix - ib(x)\)（\(x \in I\)）定义的连续线性映射。\(u_{b}\) 的图像是 u 的图像与该线性映射的图像的直和；因此它是闭的。闭算子 \(u_{b}\) 是 u 的扩张且满足 \(u_{b} \subset u^{*}\)。

设 \(\gamma_{1}\) 与 \(\eta_{1}\) 是 I 的元素。考虑元素

\[
\gamma = \gamma_ {1} + b (\gamma_ {1}), \quad \eta = \eta_ {1} + b (\eta_ {1})
\]

它们是图像 \(\Gamma_{b|I}\) 中的元素。我们计算

\[
\begin{array}{c} \langle \gamma | u ^ {*} (\eta) \rangle = i \big (\langle b (\gamma_ {1}) | \eta_ {1} \rangle - \langle \gamma_ {1} | b (\eta_ {1}) \rangle \big) + i \big (\langle \gamma_ {1} | \eta_ {1} \rangle - \langle b (\gamma_ {1}) | b (\eta_ {1}) \rangle \big) \\ = i \big (\langle b (\gamma_ {1}) | \eta_ {1} \rangle - \langle \gamma_ {1} | b (\eta_ {1}) \rangle \big) \end{array}
\]

因为 b 在 I 上等距。由此我们推出

\[
\langle u ^ {*} (\gamma) \mid \eta \rangle = \langle \gamma \mid u ^ {*} (\eta) \rangle\tag{12}
\]

对一切 \(\gamma\) 与 \(\eta\)（在 \(\Gamma_{b|I}\) 中）成立。

设 x 与 y 是 \(\operatorname{dom}(u)\) 的元素，\(\gamma\) 与 \(\eta\) 是 \(\Gamma_{b|I}\) 的元素。我们得到

\[
\begin{array}{c} \langle x + \gamma | u _ {b} (y + \eta) \rangle = \langle x | u (y) \rangle + \langle x | u ^ {*} (\eta) \rangle + \langle \gamma | u (y) \rangle + \langle \gamma | u ^ {*} (\eta) \rangle \\ = \langle u (x) | y \rangle + \langle u (x) | \eta \rangle + \langle u ^ {*} (\gamma) | y \rangle + \langle u ^ {*} (\gamma) | \eta \rangle \\ = \langle u _ {b} (x + \gamma) | y + \eta \rangle , \end{array}
\]

其中用到了 u 对称这一事实、\(\gamma\) 与 \(\eta\) 属于 \(u^{*}\) 的定义域这一事实，以及公式 (12)。于是算子 \(u_{b}\) 是 u 的闭对称扩张。断言 a) 得证。

映射 \(b \mapsto u_b\) 是单射。事实上，从 \(E_+\) 到 \(E_-\) 中的部分等距映射由其初始子空间及其在该子空间上的限制唯一确定。而 \(u_b\) 的定义域确定了 \(b\) 的初始子空间 I 以及 \(b\) 在 I 上限制的图像。

我们来证明映射 \(b \mapsto u_b\) 是满射。设 \(w\) 是 \(u\) 的闭对称扩张。我们有 \(w \subset w^* \subset u^*\)。于是 \(w\) 的定义域是 \(\text{dom}(u^*)\) 中包含 \(\text{dom}(u)\) 的子空间，且 \(w\) 是 \(u^*\) 到该子空间上的缩减。设 G 是 \(w\) 的定义域与 \(E_+ \oplus E_-\) 的交。它是 \(E_{u^*}\) 的闭子空间（引理 7），且由命题 26 有 \(\text{dom}(w) = \text{dom}(u) \oplus G\)。

设 \(x \in \mathrm{E}_{+}\)，\(y \in \mathrm{E}_{-}\) 满足 \(x + y \in \mathrm{G}\)。由于 \(w\) 对称且是 \(u^{*}\) 到 \(\operatorname{dom}(w)\) 上的缩减，我们有

\[
\langle x + y \mid u ^ {*} (x + y) \rangle = \langle u ^ {*} (x + y) \mid x + y \rangle .
\]

此等式等价于

\[
\begin{array}{c} \langle x | i x \rangle + \langle x | - i y \rangle + \langle y | i x \rangle + \langle y | - i y \rangle \\ = \langle i x | x \rangle + \langle i x | y \rangle + \langle - i y | x \rangle + \langle - i y | y \rangle , \end{array}
\]

亦即等价于 \(\|x\|^{2} = \|y\|^{2}\)。特别地，典范投影 \(p_{+}: G \to E_{+}\) 是单射，且若用 I 表示它的像，则空间 G 是等距映射 \(b_{0}\)（从 I 到 \(E_{-}\) 中的）的图像。对 I 中的 x 与 \(E_{-}\) 中满足 \(x + y \in G\) 的 y，我们有

\[
\begin{array}{r} \| x + y \| _ {u ^ {*}} ^ {2} = \| x \| _ {u ^ {*}} ^ {2} + \| y \| _ {u ^ {*}} ^ {2} = 2 (\| x \| ^ {2} + \| y \| ^ {2}) \\ = 4 \| x \| ^ {2} = 4 \| p _ {+} (x + y) \| ^ {2}. \end{array}
\]

于是，把 G 看作希尔伯特空间 \(E_{u^{*}}\) 的闭子空间时，映射 \(x \mapsto \frac{1}{2}p_{+}(x)\) 是从 G 到 I 上的等距映射。特别地，I 在 \(E_{+}\) 中闭（I，第 107 页，引理 8）。于是存在 u 的唯一的边界条件 b，以 I 为初始空间且在 I 上与 \(b_{0}\) 相合。我们有 \(\text{dom}(w) = \text{dom}(u_{b})\)，从而 \(w = u_{b}\)。这就建立了断言 b)。

我们来证明断言 \(c\)）与 \(d\)）。设 \(b\) 是 \(u\) 的一个边界条件，I 是它的初始子空间。对 \(x_0 \in \mathrm{dom}(u)\) 与 \(y \in \mathrm{I}\)，有

\[
u _ {b} (x _ {0} + y + b (y)) - i (x _ {0} + y + b (y)) = (u - i) (x _ {0}) - 2 i b (y).
\]

我们有 \(\operatorname{Im}(u - i) = \operatorname{E}_{-}^{\circ}\)（命题 25）。由于 \(b\) 的像包含在 \(\operatorname{E}_{-}\) 中，该公式表明 \(\operatorname{Im}(u_b - i) = \operatorname{Im}(u - i) \oplus \operatorname{Im}(b)\)。由同上所引，我们因此有 \(\operatorname{Ker}(u_b^* + i) = \operatorname{E}_{-} \cap \operatorname{Im}(b)^{\circ} = \operatorname{Ker}(b^*)\)。

类似地，可以验证 \(\operatorname{Im}(u_b + i) = \operatorname{Im}(u + i) \oplus \operatorname{I}\)，从而（同上所引）有 \(\operatorname{Ker}(u_b^* - i) = \operatorname{I}^\circ \cap \operatorname{E}_+ = \operatorname{Ker}(b)\)。

推论 1. --- 设 u 是复希尔伯特空间 E 上的闭对称算子。映射 \(b \mapsto u_b\) 诱导出从 \(\text{Ker}(u^* - i)\) 到 \(\text{Ker}(u^* + i)\) 上的等距同构集合与 u 的自伴扩张集合之间的双射。特别地，存在 u 的自伴扩张，当且仅当 u 的亏指标的两个分量相等。

这由命题 26 的推论与该定理以及 EVT，V，第 25 页，推论 2 得出。

推论 2. --- 设 E 是实希尔伯特空间。设 u 是 E 上的闭对称算子。部分算子 \(u_{(\mathbf{C})}\)（\(\mathrm{E}_{(\mathbf{C})}\) 上的）对称且闭，并且至少有一个自伴扩张。

部分算子 \(u_{(\mathbf{C})}\) 闭且对称（IV，第 228 页命题 2 与 IV，第 239 页引理 3）。同构 \(j\)（从 \(\mathrm{E}_{(\mathbf{C})}\) 到 \(\mathrm{E}_{(\mathbf{C})}\) 上的 R-线性同构）满足 \(j(z \otimes x) = \overline{z} \otimes x\)，对一切 \(z \in \mathbf{C}\) 与每个 \(x \in \mathrm{E}\)；它诱导出从 \(\mathrm{Ker}(u_{(\mathbf{C})}^* - i)\) 到 \(\mathrm{Ker}(u_{(\mathbf{C})}^* + i)\) 上的希尔伯特空间同构，于是可以应用推论 1。

推论 3. --- 设 u 是 E 上的对称算子。下列条件等价：

\begin{enumerate}
\def\labelenumi{(\roman{enumi})}
\item
  算子 u 本质自伴；
\item
  空间 \(\operatorname{Ker}(u^{*} + i)\) 与 \(\operatorname{Ker}(u^{*} - i)\) 为零；
\item
  空间 \(\operatorname{Im}(u + i)\) 与 \(\operatorname{Im}(u - i)\) 在 E 中稠密；
\item
  对称部分算子 u 有唯一的自伴扩张。
\end{enumerate}

由 IV，第 236 页命题 7，c)，条件 (ii) 与 (iii) 等价。

部分算子 \(u\) 本质自伴，当且仅当 \(\overline{u}\) 自伴，亦即当且仅当 \(\text{Ker}(\overline{u}^* - i) = \text{Ker}(\overline{u}^* + i) = \{0\}\)（命题 26 的推论）。由于 \(\overline{u}^* = u^*\)（IV，第 237 页，命题 8），条件 (i) 因此等价于条件 (ii)。此外，前一推论表明 (ii) 蕴含 \(u\) 有唯一的自伴扩张，此即条件 (iv)。

最后，设 \(u\) 有唯一的自伴扩张 \(v\)。那么 \(\overline{u}\) 亦然，因为 \(\overline{u}\) 的每个自伴扩张都是 \(u\) 的自伴扩张，从而等于 \(v\)。由前一推论，空间 \(\mathrm{Ker}(\overline{u}^{*} + i)\) 与 \(\mathrm{Ker}(\overline{u}^{*} - i)\) 必为零，再利用 \(\overline{u}^{*} = u^{*}\) 即得条件 (ii)。

例。--- 设 U 是 \(\mathbf{R}^{n}\) 中赋有勒贝格测度 \(\mu\) 的开子集。设 \(\Delta\) 是数量微分算子

\[
\Delta = - \sum_ {i = 1} ^ {n} \partial_ {i} ^ {2}
\]

此算子定义在 U 上。用 \(u\) 表示实希尔伯特空间 \(\mathrm{L}_{\mathbf{R}}^{2}(\mathrm{U},\mu)\) 上以 \(\mathcal{D}_{\mathbf{R}}(\mathrm{U})\)（IV，第 201 页，注）为定义域、由 \(\varphi \mapsto -\sum_{i=1}^{n}\partial_{i}^{2}\varphi\) 定义的部分算子。我们有 \(u_{(\mathbf{C})} = \Delta_{-}\)，它是可闭的部分算子（IV，第 242 页，命题 13）且对称（因为 \(\Delta\) 形式上对称），故 \(u\) 可闭且对称（IV，第 228 页命题 2 与 IV，第 239 页引理 3）。于是由上述推论，存在 \(u_{(\mathbf{C})}\) 的自伴扩张。这就是 U 上的一个拉普拉斯算子（IV，第 243 页）。

\section{正规部分算子与谱定理}

\subsection{有界化}

设 E 是复希尔伯特空间。我们用 s 表示酉映射 \((x,y)\mapsto(-y,x)\)（\(E\oplus E\) 上的）。设 u 是 E 上具有稠密定义域的闭部分算子。部分算子 \(u^{*}u\) 自伴且正（IV，第 241 页，命题 12），故 \(-1\notin\mathrm{Sp}(u^{*}u)\)（IV，第 248 页，命题 17）。我们记 \(\mathrm{W}(u)=(1_{\mathrm{E}}+u^{*}u)^{-1}=-\mathrm{R}(u^{*}u,-1)\)；它是 E 的正的单射自同态。

用 \(p_{1}\) 与 \(p_{2}\) 表示 \(\Gamma_{u}\) 到 E 中的两个典范投影。它们是 \(\mathcal{L}(\Gamma_{u};\mathrm{E})\) 的元素，且有对应关系的等式 \(p_{2}=u\circ p_{1}\)。设 \((j,|p_{1}^{*}|)\) 是自同态 \(p_{1}^{*}\in\mathcal{L}(\mathrm{E};\Gamma_{u})\) 的极分解（I，第 140 页，定义 4），即 \(p_{1}^{*}=|p_{1}^{*}|\circ j\)。映射 j 是从 E 到 \(\Gamma_{u}\) 中的部分等距。

定义 1. --- E 的自同态 \(p_2 \circ j\) 称为 \(u\) 的有界化；我们用 \(b(u)\) 表示这个自同态。

我们有 \(\|b(u)\|\leqslant1\)。

命题 1. --- 我们有下列公式

\[
| p _ {1} ^ {*} | = \mathrm{W} (u) ^ {1 / 2},\tag{1}
\]

\[
b (u) = u \circ \mathrm{W} (u) ^ {1 / 2} = u \circ | p _ {1} ^ {*} |,\tag{2}
\]

\[
1 _ {\mathrm{E}} - b (u) ^ {*} b (u) = \mathrm{W} (u),\tag{3}
\]

\[
b (u) \mathrm{W} (u) = \mathrm{W} (u ^ {*}) b (u), \quad b (u) \mathrm{W} (u) ^ {1 / 2} = \mathrm{W} (u ^ {*}) ^ {1 / 2} b (u),\tag{4}
\]

设 \(x \in E\)。令 \(y = \mathrm{W}(u)(x) \in \mathrm{dom}(u^{*}u)\)。我们有 \(y \in \mathrm{dom}(u)\) 且 \(u(y) \in \mathrm{dom}(u^{*})\)。对形如 \((y_{1}, u(y_{1}))\) 的、\(\Gamma_{u}\) 中的每个元素（其中 \(y_{1} \in \mathrm{dom}(u)\)），计算得

\[
\begin{array}{c} \langle (y, u (y)) | (y _ {1}, u (y _ {1})) \rangle = \langle y | y _ {1} \rangle + \langle u (y) | u (y _ {1}) \rangle \\ = \langle y + u ^ {*} u (y) | y _ {1} \rangle = \langle x | p _ {1} (y _ {1}, u (y _ {1})) \rangle . \end{array}
\]

这就表明 \(p_{1}^{*}(x) = (y, u(y))\)，从而 \(p_{1} \circ p_{1}^{*}(x) = y = \mathrm{W}(u)(x)\)。于是有 \(p_{1} \circ p_{1}^{*} = \mathrm{W}(u)\)，故 \(|p_{1}^{*}| = \mathrm{W}(u)^{1/2}\)，即公式 (1)。特别地，由 I，第 140 页命题 11 的推论有 \(\operatorname{Im}(p_{1}) = \operatorname{Im}(|p_{1}^{*}|)\)，由此得出 \(\operatorname{dom}(u) = \operatorname{Im}(p_{1}) = \operatorname{Im}(\mathrm{W}(u)^{1/2})\)。

于是部分算子 \(u \circ \mathrm{W}(u)^{1/2}\) 以 E 为定义域。关系式 \(p_1 \circ j = (p_1^*)^* \circ j = |p_1^*| = \mathrm{W}(u)^{1/2}\)（I，第 140 页，命题 11，a)）进而蕴含

\[
b (u) = p _ {2} \circ j = u \circ p _ {1} \circ j = u \circ \mathrm{W} (u) ^ {1 / 2},
\]

此即公式 (2)。

由 \(\operatorname{Ker}(j) = \operatorname{Ker}(|p_1^*|)\)（I，第 139 页，命题 10，\(b\)）与公式 (1) 便推出部分等距 \(j: \mathrm{E} \to \Gamma_u\) 是单射，从而是等距的，于是 \(j^* \circ j = 1_{\mathrm{E}}\)。把 j 写成 \(j = (p_1 \circ j, b(u))\)，我们得到

\[
1 _ {\mathrm{E}} = j ^ {*} \circ j = (p _ {1} \circ j) ^ {*} (p _ {1} \circ j) + b (u) ^ {*} b (u) = \mathrm{W} (u) + b (u) ^ {*} b (u),
\]

由此得公式 (3)。

设 \(x \in \mathrm{dom}(u)\)，令 \(y = \mathrm{W}(u)(x)\)。我们有 \(y \in \mathrm{dom}(u)\)，且公式 \(y + u^{*}u(y) = x\) 蕴含 \(u^{*}u(y) \in \mathrm{dom}(u)\)。于是我们有

\[
\begin{array}{r l} & {\mathrm{W} (u ^ {*}) ^ {- 1} (u (y)) = (1 _ {\mathrm{E}} + u u ^ {*}) (u (y))} \\ & {\qquad = u (y) + u u ^ {*} u (y) = u (y + u ^ {*} u (y)) = u (x),} \end{array}
\]

这表明 \(u \circ \mathrm{W}(u)\) 到 u 的定义域上的缩减等于 \(\mathrm{W}(u^{*}) \circ u\)。由于 \(|p_{1}^{*}|\) 的像等于 \(p_{1}\) 的像，从而等于 u 的定义域（I，第 140 页命题 11 的推论），由此得

\[
u \circ \mathrm{W} (u) \circ | p _ {1} ^ {*} | = \mathrm{W} (u ^ {*}) \circ u \circ | p _ {1} ^ {*} | = \mathrm{W} (u ^ {*}) \circ b (u)
\]

（公式 (2)）。此外，\(|p_{1}^{*}| = \mathrm{W}(u)^{1/2}\) 与 \(\mathrm{W}(u)\) 交换，于是我们得到

\[
b (u) \circ \mathrm{W} (u) = u \circ | p _ {1} ^ {*} | \circ \mathrm{W} (u) = u \circ \mathrm{W} (u) \circ | p _ {1} ^ {*} | = \mathrm{W} (u ^ {*}) \circ b (u).
\]

此关系式蕴含 \(b(u) \circ f(\mathrm{W}(u)) = f(\mathrm{W}(u^{*})) \circ b(u)\)，对每个函数 \(f \in \mathcal{C}(\mathbf{R}_{+})\)（I，第 113 页，命题 11），特别地有 \(b(u) \circ \mathrm{W}(u)^{1/2} = \mathrm{W}(u^{*})^{1/2} \circ b(u)\)。这就证明了公式 (4) 并完结了证明。

推论。--- 我们有 \(b(u)^{*} = b(u^{*})\)。

用 \(q_{1}\) 与 \(q_{2}\) 表示投影 \(\Gamma_{u^{*}}\to\mathrm{E}\)，并设 \((k,|q_{1}^{*}|)\) 是 \(q_{1}^{*}\) 的极分解，于是 \(b(u^{*})=q_{2}\circ k\)。对 E 中所有 \(x\) 与 \(y\)，我们有 \(j(x)\in\Gamma_{u}\) 与 \(k(y)\in\Gamma_{u^{*}}\)，且由 IV，第 236 页命题 7，\(\Gamma_{u}\) 与 \(s(\Gamma_{u^{*}})\) 正交，由此得

\[
0 = \langle j (x) | s (k (y)) \rangle = \langle p _ {1} \circ j (x) | - q _ {2} \circ k (y) \rangle + \langle p _ {2} \circ j (x) | q _ {1} \circ k (y) \rangle .
\]

由于 \(p_1 \circ j = |p_1^*|\) 且 \(q_1 \circ k = |q_1^*|\)（I，第 140 页，命题 11，b)），我们得到

\[
\langle b (u) x \mid | q _ {1} ^ {*} | y \rangle = \langle | p _ {1} ^ {*} | x \mid b (u ^ {*}) y \rangle ,
\]

故 \(|q_1^*|b(u) = b(u^*)^* |p_1^*|\)。利用公式 (1) 与 (4)，我们由此得出 \(b(u)|p_1^*| = |q_1^*|b(u) = b(u^*)^* |p_1^*|\)，又因 \(|p_1^*|\) 的像是 \(u\) 的定义域，它在 E 中稠密，故得 \(b(u) = b(u^*)^*\)。

我们用 \(\Omega(\mathrm{E})\) 表示由 \(v \in \mathcal{L}(\mathrm{E})\)（满足 \(1_{\mathrm{E}} - v^{*}v\) 正且单射的）所成的集合。对 \(v \in \Omega(\mathrm{E})\)，埃尔米特自同态 \((1_{\mathrm{E}} - v^{*}v)^{1/2}\) 是单射的，因为它的平方是单射的；我们用 \(\mathrm{B}(v)\) 表示部分算子 \(v \circ ((1_{\mathrm{E}} - v^{*}v)^{1/2})^{-1}\)。

注意 \(1_{\mathrm{E}} - v^{*}v\) 正当且仅当 \(\| v\| \leqslant 1\)，因为 \(\langle x|(1_{\mathrm{E}} - v^{*}v)(x)\rangle = \| x\| ^2 -\| v(x)\| ^2\) 对一切 \(x\in \mathrm{E}\) 成立。

引理 1. --- 子集 \(\Omega(\mathrm{E})\) 在 \(\mathcal{L}(\mathrm{E})\) 中是自伴的。

设 \(v \in \Omega(\mathrm{E})\)。我们有 \(\|v^{*}\| = \|v\| \leqslant 1\)。于是自同态 \(1_{\mathrm{E}} - vv^{*}\) 是正的。它是单射的：若 \(x \in \operatorname{Ker}(1_{\mathrm{E}} - vv^{*})\)，我们有 \(vv^{*}(x) = x\)，故 \(v^{*}(v(v^{*}(x))) = v^{*}(x)\)，于是 \(v^{*}(x) = 0\)（因为 \(1_{\mathrm{E}} - v^{*}v\) 单射），最后 \(x = v(v^{*}(x)) = 0\)。引理由此得证。

命题 2. — 映射 \(u \mapsto b(u)\) 定义了从 E 上闭的稠定部分算子集合到集合 \(\Omega(E)\) 上的一个双射。逆双射由 \(v \mapsto B(v)\) 给出。

关系式 \(1_{\mathrm{E}} - b(u)^{*}b(u) = \mathrm{W}(u)\)（公式 (3)）蕴含 \(b(u)\) 属于 \(\Omega (\mathrm{E})\)，因为 \(\mathrm{W}(u)\) 正且单射。此外，由 \(b(u) = u\circ \mathrm{W}(u)^{1 / 2}\)（公式 (2)），我们得到

\[
u = b (u) \circ (\mathrm{W} (u) ^ {1 / 2}) ^ {- 1} = b (u) \circ ((1 - b (u) ^ {*} b (u)) ^ {1 / 2}) ^ {- 1} = \mathrm{B} (b (u)).
\]

反之，设 \(v \in \Omega(\mathrm{E})\)。令 \(w = (1_{\mathrm{E}} - v^{*}v)^{1/2}\)，并令 \(u = \mathrm{B}(v) = v \circ w^{-1}\)。u 的定义域是 w 的像，由于 w 埃尔米特且单射，它在 E 中稠密（EVT，V，第 41 页，命题 4 (i)）。于是对每个 \(x \in \mathrm{dom}(u)\)，有

\[
\begin{array}{r l} & {\| u (x) \| ^ {2} = \langle (v ^ {*} v) (w ^ {- 1} (x)) | w ^ {- 1} (x) \rangle} \\ & {\qquad = - \langle (1 _ {\mathrm{E}} - v ^ {*} v) (w ^ {- 1} (x)) | w ^ {- 1} (x) \rangle + \| w ^ {- 1} (x) \| ^ {2}} \\ & {\qquad = - \langle w (x) | w ^ {- 1} (x) \rangle + \| w ^ {- 1} (x) \| ^ {2}} \end{array}
\]

因为 \((1_{\mathrm{E}} - v^{*}v) \circ w^{-1}\) 是 w 到 \(w^{-1}\) 的定义域上的缩减。由于 w 自伴，由此得

\[
\| w ^ {- 1} (x) \| ^ {2} = \| x \| ^ {2} + \| v \circ w ^ {- 1} (x) \| ^ {2} = \| x \| ^ {2} + \| u (x) \| ^ {2},
\]

于是由 IV，第 231 页引理 2，部分算子 \(\mathrm{B}(v) = u = v \circ w^{-1}\) 是闭的。

最后我们来证明 \(b(\mathrm{B}(v)) = v\)。由于 \(u = v \circ w^{-1}\)，我们有 \(v = u \circ w\)，因此由公式 (2)，只需验证 \(w = \mathrm{W}(u)^{1/2}\)，甚至只需验证 \(1_{\mathrm{E}} - v^{*}v = \mathrm{W}(u)\)。

自同态 \(v^*\) 属于 \(\Omega(\mathrm{E})\)。由引理 1，令 \(w' = (1_{\mathrm{E}} - vv^{*})^{1/2}\)。它是 E 的单射正自同态，且我们有 \(v \circ w = w' \circ v\)（I，第 113 页，命题 11）。\(u\) 的图像是形如 \((w(x), v(x))\)（\(x \in \mathrm{E}\)）的元素所成的集合。由 IV，第 236 页命题 7，\(u^*\) 的图像是 \(s(\Gamma_u)^{\circ}\)。由于对一切 \(x\) 与 \(y \in \mathrm{E}\)，有

\[
\langle (w ^ {\prime} (x), v ^ {*} (x)) | (- v (y), w (y)) \rangle = - \langle x | w ^ {\prime} v (y) \rangle + \langle x | v w (y) \rangle = 0,
\]

故 \(u^*\) 的图像包含形如 \((w'(x), v^*(x))\)（\(x \in \mathrm{E}\)）的元素，并且此时有 \(u^*(w'(x)) = v^*(x)\)。特别地，\(u^*\) 的定义域包含 \(w'\) 的像。

设 \(x \in E\)，令 \(y = w^{2}(x) = (1_{\mathrm{E}} - v^{*}v)(x)\)，于是 \(x = y + v^{*}v(x)\)。我们有 \(y \in \operatorname{dom}(u)\) 且 \(u(y) = v(w^{-1}(y)) = v(w(x)) = w'(v(x))\)。特别地，\(u(y) \in \operatorname{Im}(w) \subset \operatorname{dom}(u^{*})\)，且 \(u^{*}(u(y)) = v^{*}(v(x))\)。由此得 \(x = y + v^{*}v(x) = y + u^{*}u(y)\)，从而 \(y = \operatorname{W}(u)(x)\)，即

\[
(1 _ {\mathrm{E}} - v ^ {*} v) (x) = \mathrm{W} (u) (x).
\]

我们于是证明了 \(1_{\mathrm{E}} - v^{*}v = \mathrm{W}(u)\)，这正是所要的。

推论。--- 自同态 \(b(u)\) 是埃尔米特的，当且仅当 u 自伴。

这由映射 \(u \mapsto b(u)\) 的单射性（命题 2）以及公式 \(b(u^{*}) = b(u)^{*}\)（命题 1 的推论）得出。

\subsection{正规部分算子与谱定理}

在本节中，E 是复希尔伯特空间。

定义 2. --- 设 u 是 E 上的部分算子。称 u 为正规的，如果 u 闭且稠定，并且它的有界变换 b(u) 是 E 的正规自同态。

若 \(u \in \mathcal{L}(\mathrm{E})\)，则由公式 \(1_{\mathrm{E}} - b(u)^{*}b(u) = \mathrm{W}(u) = (1_{\mathrm{E}} + u^{*}u)^{-1}\)（IV，第 262 页，命题 1）与 \(b(u^{*}) = b(u)^{*}\)（同上所引之推论），此定义与 EVT，V，第 42 页，定义 4 相一致。

若 \(u\) 是 \(\mathrm{E}\) 上的自伴部分算子，则 \(b(u)\) 是埃尔米特的（IV，第 264 页命题 2 的推论），故 \(u\) 正规。

设 D 是 C 中的开单位圆盘。我们用 \(\beta\) 表示由 \(\beta(z)=z/\sqrt{1+|z|^{2}}\) 定义的从 C 到 D 中的函数。它是一个同胚，其逆满足 \(\beta^{-1}(z)=z/\sqrt{1-|z|^{2}}\)，对 \(z\in D\)。

设 \(u \in \mathcal{L}(\mathrm{E})\)。由 IV，第 262 页命题 1 的公式 (2) 与 (3) 得 \(u = \beta^{-1}(b(u))\)，从而

\[
b (u) = \beta (u).\tag{5}
\]

引理 2. --- 设 X 是局部紧拓扑空间，μ 是 X 上的正测度。

\begin{enumerate}
\def\labelenumi{\alph{enumi})}
\item
  设 \(g\) 是 X 上 \(\mu\)-可测的函数。乘法算子 \(m_g\)（\(\mathrm{L}^2 (\mathrm{X},\mu)\) 上的）是正规的，且 \(b(m_g) = m_{\beta \circ g}\)；
\item
  设 \(h\colon X\to D\) 是 \(\mu\)-可测函数。自同态 \(m_h\) 属于 \(\Omega(L^2(X,\mu))\)，且 \(B(m_h)=m_{\beta^{-1}\circ h}\)。
\end{enumerate}

部分算子 \(m_g\) 闭且稠定。由于 \(m_g^* m_g = m_{\overline{g}} m_g = m_{|g|^2}\)（IV，第 253 页命题 23 与 IV，第 255 页命题 24），我们有

\[
\mathrm{W} (m _ {g}) = - \mathrm{R} (m _ {| g | ^ {2}}, - 1) = m _ {(1 + | g | ^ {2}) ^ {- 1}}
\]

（IV，第 252 页，命题 22）。由此得

\[
b (m _ {g}) = m _ {g} \circ \mathrm{W} (m _ {g}) ^ {1 / 2} = m _ {g} \circ m _ {(1 + | g | ^ {2}) ^ {- 1 / 2}} = m _ {\beta \circ g}
\]

这里依次用到 IV，第 262 页的公式 (2)、IV，第 187 页命题 6 的推论以及 IV，第 254 页引理 6。最后应用 IV，第 187 页命题 6，即得断言 a)。

我们来证明 \(b\)）。由于 \(1 - |h(x)|^2 > 0\) 对一切 \(x \in \mathrm{X}\) 成立，我们有 \(m_h \in \Omega(\mathrm{L}^2(\mathrm{X}, \mu))\)（IV，第 187 页命题 6 的推论）。由于映射 \(u \mapsto b(u)\) 是单射（IV，第 264 页，命题 2），且自同态 \(b(\mathrm{B}(m_h))\) 与 \(b(m_{\beta^{-1}\circ h})\) 都等于 \(m_h\)（由断言 \(a\)），故确有 \(\mathrm{B}(m_h) = m_{\beta^{-1}\circ h}\)。

定理 1（“谱定理”）。--- 设 u 是希尔伯特空间 E 上的正规部分算子。存在局部紧拓扑空间 X、X 上的正测度 μ、从 L²(X,μ) 到 E 上的等距同构 θ 以及 X 上的连续函数 g，使得 \(u = \theta \circ m_g \circ \theta^{-1}\)。

按定义，E 的自同态 \(b(u)\) 是正规的。由 E 的正规自同态的谱定理（IV，第 193 页，推论 1），存在局部紧拓扑空间 \(\widetilde{X}\)、\(\mu\)（\(\widetilde{X}\) 上的正测度）、\(\widetilde{\theta}\)（从 \(\mathrm{L}^2 (\widetilde{\mathrm{X}},\mu)\) 到 E 上的等距同构），以及 \(h\)（\(\widetilde{X}\) 上的有界连续函数），使 \(b(u)\) 与 \(\widetilde{\theta}\circ m_h\circ \widetilde{\theta}^{-1}\) 相合。

设 N 是由 \(x \in \widetilde{X}\)（满足 \(|h(x)| \geqslant 1\) 的）所成的闭子集。由于自同态 \(1 - b(u)^{*}b(u) = \widetilde{\theta} \circ m_{1-|h|^2} \circ \widetilde{\theta}^{-1}\) 正且单射，集合 N 是局部 \(\mu\)-零测的（IV，第 186 页引理 7 与 IV，第 187 页命题 6 的推论）。于是令 X = \(\widetilde{X} - N\)。它是局部紧空间，且函数到 X 上的限制诱导出从 \(\mathrm{L}^{2}(\widetilde{\mathrm{X}}, \mu)\) 到 \(\mathrm{L}^{2}(\mathrm{X}, \mu|\mathrm{X})\) 上的等距同构（IV，第 179 页，命题 1）。于是 \(\widetilde{\theta}\) 与函数到 X 上的限制的复合诱导出等距同构 \(\theta\)（从 \(\mathrm{L}^{2}(\mathrm{X}, \mu|\mathrm{X})\) 到 E 上的）。由于 \(b(u) = \theta \circ m_{h|\mathrm{X}} \circ \theta^{-1}\)，得 \(u = \theta \circ m_{\beta^{-1} \circ (h|\mathrm{X})} \circ \theta^{-1}\)（引理 2）。定理由此公式得出。

引理 3. --- 设 u 是 E 上的正规部分算子。我们有

\[
\operatorname{Sp} (b (u)) \cap \mathrm{D} = \beta (\operatorname{Sp} (u)).
\]

我们有 \(b(u) = u \circ \mathrm{W}(u)^{1/2}\)（IV，第 262 页，命题 1）。自同态 \(\mathrm{W}(u)^{1/2}\) 是单射的，且它的像是 u 的定义域（IV，第 262 页，公式 (1)）。

设 \(\lambda \in \mathbf{C}\)。数 \(\lambda\) 属于 \(u\) 的预解集，当且仅当 \((u - \lambda 1_{\mathrm{E}}) \circ \mathrm{W}(u)^{1/2}\) 是从 E 到 E 上的双射线性映射（IV，第 245 页，注 3）。而由 IV，第 262 页的公式 (3)，有

\[
(u - \lambda 1 _ {\mathrm{E}}) \circ \mathrm{W} (u) ^ {1 / 2} = b (u) - \lambda (1 _ {\mathrm{E}} - b (u) ^ {*} b (u)) ^ {1 / 2} = f _ {\lambda} (b (u))
\]

其中 \(f_{\lambda}\) 是在 \(\overline{\mathrm{D}}\) 上由 \(z \mapsto z - \lambda (1 - |z|^2)^{1/2}\) 定义的连续函数。自同态 \(f_{\lambda}(b(u))\) 是双射的，当且仅当它的谱不含 0。由于 \(\mathrm{Sp}(f_{\lambda}(b(u))) = f_{\lambda}(\mathrm{Sp}(b(u)))\)（I，第 111 页，命题 7，推论 2），此事成立当且仅当 0 不属于集合 \(f_{\lambda}(\mathrm{Sp}(b(u)))\)。于是，\(\lambda \in \mathrm{Sp}(u)\) 当且仅当存在 \(z \in \mathrm{Sp}(b(u)) \cap \overline{\mathrm{D}}\) 使得

\[
z - \lambda (1 - | z | ^ {2}) ^ {1 / 2} = 0.
\]

此等式若成立，便蕴含 \(z \in D\)，且意味着 \(\lambda = \beta^{-1}(z)\)。我们得出结论 \(\operatorname{Sp}(u) = \beta^{-1}(\operatorname{Sp}(b(u)) \cap D)\)，正如所述。

设 u 是 E 上的正规部分算子。设 \(f \in \mathcal{K}(\mathrm{Sp}(u))\)。映射 \(z \mapsto f(\beta^{-1}(z))\)（从 \(\mathrm{Sp}(b(u)) \cap D\) 到 C 中的）连续且具有紧支集；于是唯一的映射 \(f_{\beta}\)（从 \(\mathrm{Sp}(b(u))\) 到 C 中、用零延拓它的）是连续的。

定义 3. --- 对每个函数 \(f \in \mathcal{K}(\mathrm{Sp}(u))\)，我们定义 E 的自同态 \(f(u)\) 如下：\(f(u) = f_{\beta}(b(u))\)。

映射 \(f \mapsto f_{\beta}\) 是从 \(\mathcal{K}(\mathrm{Sp}(u))\) 到 \(\mathcal{C}(\mathrm{Sp}(b(u)))\) 中的复代数同态，故映射 \(f \mapsto f(u)\) 是从 \(\mathcal{K}(\mathrm{Sp}(u))\) 到 \(\mathcal{L}(\mathrm{E})\) 中的复代数同态。我们有 \(\overline{f}(u) = f(u)^*\)，因为 \(\overline{f}_{\beta} = \overline{f_{\beta}}\)。若 \(f \geqslant 0\)，则 \(f_{\beta} \geqslant 0\)，故 \(f(u)\) 是 E 的正自同态。

注记。--- 设 u 是 E 的正规自同态，于是 \(\mathrm{Sp}(u)\) 是 C 的紧子集。由于有界变换 \(b(u)\) 与 \(\beta(u)\) 相合（公式 (5)），且 \(\mathrm{Sp}(b(u)) = \beta(\mathrm{Sp}(u))\)（I，第 111 页，命题 7，推论 2），\(b(u)\) 的谱是 D 的紧子集。对每个连续函数 \(f \in \mathcal{C}(\mathrm{Sp}(u))\)，函数 \(f_{\beta}\) 与 \(f \circ \beta^{-1}\) 到 \(\mathrm{Sp}(u)\) 上的限制相合。于是，\(f_{\beta}(\beta(u))\) 与由 u 的连续泛函演算所定义的自同态 \(f(u)\) 相合。因此上述定义与星代数 \(\mathcal{L}(\mathrm{E})\) 的连续泛函演算的定义相容。

设 \(f\in \mathcal{K}(\mathrm{Sp}(u))\)。我们有

\[
\| f (u) \| \leqslant \| f _ {\beta} \| _ {\infty} = \| f \| _ {\infty},
\]

于是，对所有 \(x\) 与 \(y\)（E 中的），我们得到估计 \(|\langle x | f(u)y\rangle| \leqslant \|x\|\|y\|\|f\|_{\infty}\)。故映射 \(f \mapsto \langle x | f(u)y\rangle\) 是 \(\mathrm{Sp}(u)\) 上总质量 \(\leqslant \|x\|\|y\|\) 的有界测度（INT，IV，第 154 页，§ 4，第 7 目）。若 \(x = y\)，它是正测度，因为 \(f(u)\) 当 \(f \geqslant 0\) 时是正的。

定义 4. --- 设 u 是复希尔伯特空间 E 上的正规部分算子。设 x 与 y 在 E 中。由 \(f \mapsto \langle x | f(u)y \rangle\)（\(f \in \mathcal{K}(\mathrm{Sp}(u))\)）定义的 Sp(u) 上的有界测度称为 \((x,y)\) 关于 \(u\) 的谱测度。当 \(x = y\) 时，我们称它为 \(x\) 关于 \(u\) 的谱测度。

注记。--- 1) 当 u 连续时，此定义与 IV，第 190 页定义 2 相一致。

\begin{enumerate}
\def\labelenumi{\arabic{enumi})}
\setcounter{enumi}{1}
\tightlist
\item
  设 \(j\) 是 \(\mathrm{Sp}(u)\) 到 \(\mathbf{C}\) 中的典范单射。由于 \(\mathrm{Sp}(u)\) 是闭的，映射 \(j\) 是 \(\mu\)-逆紧的，对每个有界测度 \(\mu\)（\(\mathrm{Sp}(u)\) 上的）（INT，V，第 68 页，§ 6，第 1 目，命题 1）。人们常把 \(u\) 的谱测度与 \(\mathbf{C}\) 上作为它们在 \(j\) 下的像的测度等同起来（参见 INT，V，第 84 页，§ 7，第 2 目）。
\end{enumerate}

从 E×E 到巴拿赫空间 \(\mathcal{M}^{1}(\mathrm{Sp}(u))\) 中把 \((x,y)\) 映为 \((x,y)\) 关于 u 的谱测度的映射是半双线性的，并且是范数 \(\leqslant 1\) 的连续映射。

引理 4. --- 设 X 是局部紧拓扑空间，\(\mu\) 是 X 上的正测度，\(g\) 是 X 上 \(\mu\)-可测的函数。设 \(m_g\) 是乘以 \(g\) 的乘法算子（\(\mathrm{L}^2 (\mathrm{X},\mu)\) 上的）。

\begin{enumerate}
\def\labelenumi{\alph{enumi})}
\item
  映射 \(f\mapsto f(m_{g})\)（从 \(\mathcal{K}(\mathrm{Sp}(m_g))\) 到 \(\mathcal{L}(\mathrm{L}^{2}(\mathrm{X},\mu))\) 中的）由 \(f\mapsto m_{f\circ g}\) 给出；
\item
  对 \(f_1\) 与 \(f_2\)（\(\mathcal{L}^2(\mathrm{X},\mu)\) 中的，其类 \(\widetilde{f}_1\) 与 \(\widetilde{f}_2\) 在 \(\mathrm{L}^2(\mathrm{X},\mu)\) 中），\((\widetilde{f}_1,\widetilde{f}_2)\) 关于 \(m_g\) 的谱测度是限制到 \(\mathrm{Sp}(m_g)\) 上的像测度 \(g(\overline{f}_1f_2\cdot \mu)\)。
\end{enumerate}

部分算子 \(m_g\) 是正规的，且 \(b(m_g) = m_{\beta \circ g} = m_{g(1+|g|^2)^{-1/2}}\)（引理 2）。此外，\(\beta(g(x))\) 属于 \(\mathrm{Sp}(m_{\beta \circ g})\)，对每个 \(x\)（位于局部 \(\mu\)-零测集 \(Y \subset X\) 之外）（IV，第 252 页命题 22 与 IV，第 181 页引理 1）。由于 \(f_\beta(\beta(g(x))) = f(g(x))\) 对一切 \(x \in X - Y\) 成立，由定义 3 与 IV，第 187 页命题 6 的推论，我们推出 \(f(m_g) = m_{f \circ g}\)。

设 \(f_1\) 与 \(f_2\) 是 \(\mathcal{L}^2(\mathrm{X},\mu)\) 中的元素，其类 \(\widetilde{f}_1\) 与 \(\widetilde{f}_2\) 在 \(\mathrm{L}^2(\mathrm{X},\mu)\) 中。由于测度 \(\nu = \overline{f}_1f_2\cdot \mu\) 有界，像测度 \(g(\nu)\) 有定义（INT，V，第 69 页，§ 6，第 1 目，注 1）。设 \(f\in \mathcal{K}(\mathrm{Sp}(m_g))\)。我们有 \(f(m_{g}) = m_{f\circ g}\)（由 \(a\)），于是

\[
\langle \widetilde {f} _ {1} \mid f (m _ {g}) \widetilde {f} _ {2} \rangle = \int_ {\mathrm{X}} \overline {{f}} _ {1} (f \circ g) f _ {2} d \mu = \int_ {\mathrm{X}} (f \circ g) (\overline {{f}} _ {1} f _ {2} d \mu) = \int_ {\mathbf {C}} f d \nu
\]

（INT，V，第 69 页，§ 6，第 1 目，公式 (1)），这就建立了断言 b)。

例。--- 设 \(n \in \mathbf{N}\)。我们给 \(\mathbf{R}^{n}\) 赋以勒贝格测度，并按 II，第 236 页推论 3 把 \(\mathbf{R}^{n}\) 与它的对偶群等同起来。用 N 表示 \(\mathbf{R}^{n}\) 上的欧几里得范数，用 \(\mathcal{F}\) 表示 \(\mathcal{S}(\mathbf{R}^{n})\) 上与 \(\mathcal{S}'(\mathbf{R}^{n})\) 上的傅里叶变换（IV，第 217 页，第 12 目）。

设 \(\Delta_{\mathcal{S}}\) 是 \(\mathrm{L}^2 (\mathbf{R}^n)\) 上以空间 \(\mathcal{S}(\mathbf{R}^n)\) 为定义域且满足

\[
\Delta_ {\mathcal {S}} (\varphi) = - \sum_ {i = 1} ^ {n} \partial_ {i} ^ {2} \varphi
\]

的部分算子，对一切 \(\varphi\in\mathcal{S}(\mathbf{R}^{n})\) 成立。于是 \(\mathcal{F}(\Delta_{\mathcal{S}}(\varphi))=4\pi^{2}\mathrm{N}^{2}\mathcal{F}(\varphi)\)，对每个函数 \(\varphi\in\mathcal{S}(\mathbf{R}^{n})\)（IV，第 220 页，注 12）。由于傅里叶变换是 \(\mathcal{S}(\mathbf{R}^{n})\) 的自同构（IV，第 219 页，命题 18），这表明部分算子 \(u=\mathcal{F}\circ\Delta_{\mathcal{S}}\circ\mathcal{F}^{-1}\) 是乘法算子 \(m_{4\pi^{2}\mathrm{N}^{2}}\) 到空间 \(\mathcal{S}(\mathbf{R}^{n})\) 上的缩减。部分算子 \(u\) 可闭且对称；它的闭包是正的自伴部分算子 \(m_{4\pi^{2}\mathrm{N}^{2}}\)（IV，第 232 页命题 6，对空间 \(\mathcal{D}(\mathbf{R}^{n})\subset\mathcal{S}(\mathbf{R}^{n})\) 应用，并用到 IV，第 202 页命题 4）。因此，部分算子 \(\Delta_{\mathcal{S}}\) 本质自伴。我们用 \(\Delta\) 表示它的闭包；它是正的自伴算子，并且是 \(\mathbf{R}^{n}\) 上唯一的拉普拉斯算子（\(cf.\) IV，第 243 页第 6 目的例以及 IV，第 257 页命题 26 的推论）。

由 IV，第 223 页命题 21，索伯列夫空间 \(\mathrm{H}^2 (\mathbf{R}^n)\) 是由 \(f\in \mathrm{L}^2 (\mathbf{R}^n)\)（使 \((1 + \mathrm{N}^2)\mathcal{F}(f)\) 属于 \(\mathrm{L}^2 (\mathbf{R}^n)\) 的）所成的集合，亦即使 \(\mathcal{F}(f)\) 属于 \(m_{4\pi^2\mathrm{N}^2}\) 的定义域。由于傅里叶变换是空间 \(\mathrm{L}^2 (\mathbf{R}^n)\) 到其自身上的等距同构（II，第 215 页，定理 1），拉普拉斯算子 \(\Delta\) 的定义域是 \(\mathrm{H}^2 (\mathbf{R}^n)\)。我们有 \(\Delta = \mathcal{F}^{-1}\circ m_{4\pi^2\mathrm{N}^2}\circ \mathcal{F}\)。特别地，\(\Delta\) 的谱等于 \(\mathbf{R}_{+}\)。